% ch1.tex — Chapter 1: Matrices, Exponential Operators and Physical Applications
% 对应原书第 1–62 页（Sec. 1.1 – 1.6.7 + Bibliography）

\chapter{Matrices, Exponential Operators and Physical Applications}
\label{ch:1}

% ============================================================
\section{Introduction}
\label{sec:1.1}

The matrices and the relevant formalism are widely exploited tools in Physics and Mathematics.
They provide the backbone of the treatment of many problems in classical and quantum Physics.
Thus we start these lectures with a short resume on the relevant properties, aimed at fixing the notation and at illustrating the technicalities we will later employ in solving mathematical problems of physical interest.

Our starting point is the equation
\begin{equation}
\hat{A}\,\underline{\zeta} = \underline{b},
\label{eq:1.1.1}
\end{equation}
where
\[
\hat{A} = \mat{\alpha_{11} & \alpha_{12} \\ \alpha_{21} & \alpha_{22}},\qquad
\underline{\zeta} = \mat{z_{1} \\ z_{2}},\qquad
\underline{b} = \mat{b_{1} \\ b_{2}},\qquad
\forall\,\alpha_{ij},z_{i},b_{i}\in\mathbb{C}.
\]
Equation~\eqref{eq:1.1.1} can be considered as the generalisation of an algebraic equation of first degree.
Provided that we accept the matrix as a generalisation of ordinary numbers, its solution can be written in the form
\begin{equation}
\underline{\zeta} = \hat{A}^{-1}\,\underline{b},
\label{eq:1.1.3}
\end{equation}
reminiscent of the solution of first-degree algebraic equations.
Such a formal solution makes sense if the matrix $\hat{A}$ is not singular, i.e.\ if it is invertible.
We have tacitly assumed that
\begin{equation}
\hat{A}\,\hat{A}^{-1} = \hat{1},
\label{eq:1.1.4}
\end{equation}
where the real-unity matrix is
\[
\hat{1} = \mat{1 & 0 \\ 0 & 1},
\]
satisfying
\[
\hat{1}\,\hat{B} = \hat{B}\,\hat{1} = \hat{B},\qquad \hat{1}\,\underline{\zeta} = \underline{\zeta}.
\]

For the two-dimensional case the inverse of $\hat{A}$ is
\[
\hat{A}^{-1} = \frac{1}{\Delta}\mat{\alpha_{22} & -\alpha_{12} \\ -\alpha_{21} & \alpha_{11}},\qquad
\Delta = \alpha_{11}\alpha_{22} - \alpha_{12}\alpha_{21}.
\]
A matrix is singular (non-invertible) if its determinant vanishes.

Along with the unity matrix we introduce the \emph{imaginary matrix}
\begin{equation}
\hat{i} = \mat{0 & 1 \\ -1 & 0},\qquad \hat{i}^{2} = -\hat{1}.
\label{eq:1.1.9}
\end{equation}

A function of a matrix is defined via its Maclaurin expansion
\[
f(\hat{A}) = \sum_{n=0}^{\infty}\frac{f^{(n)}(0)}{n!}\,\hat{A}^{n}.
\]
Exponentiating $\hat{i}$ gives the Euler formula for matrices:
\begin{align}
e^{\alpha\hat{i}}
&= \sum_{n=0}^{\infty}\frac{\alpha^{n}}{n!}\,\hat{i}^{n}
 = \cos\alpha\,\hat{1} + \sin\alpha\,\hat{i} \notag \\
&= \mat{\cos\alpha & \sin\alpha \\ -\sin\alpha & \cos\alpha}
 = \hat{R}(\alpha).
\label{eq:1.1.17}
\end{align}
Thus the exponentiation of the imaginary matrix yields the two-dimensional rotation matrix.

The \emph{hyperbolic unit matrix}
\[
\hat{h} = \mat{0 & 1 \\ 1 & 0},\qquad \hat{h}^{2} = \hat{1},
\]
gives
\[
e^{\alpha\hat{h}} = \mat{\cosh\alpha & \sinh\alpha \\ \sinh\alpha & \cosh\alpha}
= \hat{R}_{h}(\alpha).
\]

The three units do not commute:
\[
[\hat{i},\hat{h}] = 2\hat{t},\qquad
\hat{t} = \mat{1 & 0 \\ 0 & -1}.
\]
Any $2\times 2$ matrix expands as
\[
\mat{\alpha_{11} & \alpha_{12} \\ \alpha_{21} & \alpha_{22}}
= a\,\hat{1} + b\,\hat{i} + c\,\hat{h} + d\,\hat{t}.
\]

The Hermitian conjugate is
\[
\hat{A}^{\dagger} = \mat{\alpha_{11}^{*} & \alpha_{21}^{*} \\ \alpha_{12}^{*} & \alpha_{22}^{*}}.
\]
We have $\hat{i}^{\dagger} = -\hat{i}$ (anti-Hermitian) while $\hat{h},\hat{t},\hat{1}$ are Hermitian.

% ============================================================
\section{Pauli Matrices}
\label{sec:1.2}

Define the ladder operators
\[
\spp = \mat{0 & 1 \\ 0 & 0},\qquad
\smm = \mat{0 & 0 \\ 1 & 0},\qquad
\spp^{2} = \smm^{2} = \hat{0},
\]
so that
\[
\hat{i} = \spp - \smm.
\]
The Pauli--$z$ matrix is
\[
\sigz = \frac{1}{2}\hat{t} = \mat{\tfrac{1}{2} & 0 \\ 0 & -\tfrac{1}{2}}.
\]

The commutation relations are
\begin{equation}
[\spp,\smm] = 2\sigz,\qquad
[\sigz,\spp] = \spp,\qquad
[\sigz,\smm] = -\smm,
\label{eq:1.2.8}
\end{equation}
which are the $SU(2)$ algebra. The standard Pauli vector is
\[
\sigx = \frac{1}{2}(\spp+\smm),\quad
\sigy = \frac{1}{2i}(\spp-\smm),\quad
\sigz = \frac{1}{2}\mat{1 & 0 \\ 0 & -1},
\]
with
\[
[\sig_{j},\sig_{k}] = i\,\varepsilon_{jkl}\,\sig_{l}.
\]

For a Hamiltonian $\hat{H} = \hbar\,\vec{\Omega}\!\cdot\!\vec{\sigma}$ the evolution operator is
\[
\hat{U}(t) = e^{-i\Omega t\,\vec{n}\!\cdot\!\vec{\sigma}}
= \cos\!\left(\frac{\Omega t}{2}\right)\hat{1}
- 2i\sin\!\left(\frac{\Omega t}{2}\right)\vec{n}\!\cdot\!\vec{\sigma}.
\]

A useful disentanglement formula for a traceless anti-Hermitian $2\times 2$ matrix is
\[
e^{\alpha(\spp-\smm)}
= e^{2h(\alpha)\sigz}\,e^{f(\alpha)\spp}\,e^{-g(\alpha)\smm},
\]
with $h(0)=f(0)=g(0)=0$ and (for this special case)
\[
e^{h(\alpha)} = \sec\alpha,\qquad
f(\alpha) = \tfrac{1}{2}\sin 2\alpha,\qquad
g(\alpha) = \tan\alpha.
\]

For a generic $2\times 2$ matrix $\hat{\alpha}$ with $\tr\hat{\alpha}\neq 0$,
\[
e^{\hat{\alpha}} = \mat{A_{11} & A_{12} \\ A_{21} & A_{22}},
\]
where (with $\Delta = (\alpha_{11}-\alpha_{22})^{2}+4\alpha_{12}\alpha_{21}$)
\begin{align*}
A_{11} &= \left\{\frac{\alpha_{11}-\alpha_{22}}{\sqrt{\Delta}}\sinh\frac{\sqrt{\Delta}}{2}
           + \cosh\frac{\sqrt{\Delta}}{2}\right\}e^{\frac{1}{2}\tr\hat{\alpha}},\\
A_{22} &= \left\{-\frac{\alpha_{11}-\alpha_{22}}{\sqrt{\Delta}}\sinh\frac{\sqrt{\Delta}}{2}
           + \cosh\frac{\sqrt{\Delta}}{2}\right\}e^{\frac{1}{2}\tr\hat{\alpha}},\\
\frac{A_{12}}{\alpha_{12}} &= \frac{A_{21}}{\alpha_{21}}
 = \frac{2}{\sqrt{\Delta}}\sinh\frac{\sqrt{\Delta}}{2}\,e^{\frac{1}{2}\tr\hat{\alpha}}.
\end{align*}

% ============================================================
\section{Applications of $2\times 2$ Matrices}
\label{sec:1.3}

\subsection{Classical Optics: Ray Propagation and ABCD Law}

A paraxial ray is described by $\mat{r & \vartheta}^{T}$.
Free space of length $d$:
\[
\hat{O} = \mat{1 & d \\ 0 & 1} = e^{d\spp}.
\]
Thin lens of focal length $f$:
\[
\hat{F} = \mat{1 & 0 \\ -1/f & 1} = e^{-\frac{1}{f}\smm},\qquad
\hat{D} = \mat{1 & 0 \\ 1/f & 1} = e^{\frac{1}{f}\smm}.
\]
Beam expander (squeeze):
\[
e^{b\sigz} = \mat{e^{b} & 0 \\ 0 & e^{-b}}.
\]
The ABCD law is
\[
\mat{r \\ \vartheta}_{\mathrm{out}}
= \mat{A & B \\ C & D}
  \mat{r \\ \vartheta}_{\mathrm{in}}.
\]

\subsection{Quantum Mechanics}

Pauli equation for spin-$\tfrac{1}{2}$ in a magnetic field $\vec{B}=B\hat{y}$:
\[
i\hbar\,\partial_{t}\underline{\Phi}
= -\frac{q\hbar}{2m}\,\vec{\sigma}\!\cdot\!\vec{B}\,\underline{\Phi}
\;\Longrightarrow\;
\partial_{t}\underline{\Phi}
= \omega_{c}(\spp-\smm)\underline{\Phi},\qquad
\omega_{c} = \frac{qB}{4m}.
\]
Solution: $\underline{\Phi}(t) = \hat{R}(\omega_{c}t)\,\underline{\Phi}_{0}$.

Two-level system:
\[
i\hbar\,\partial_{t}\mat{a_{+} \\ a_{-}}
= \bigl(\omega\sigz + 2\Omega\sigy\bigr)\mat{a_{+} \\ a_{-}},
\]
with the $SU(2)$ rotation evolution matrix given above.

\subsection{Particle Physics: Kaon Mixing}

Neutral kaon states:
\[
|K^{0}\rangle = |\bar{s}d\rangle,\qquad
|\bar{K}^{0}\rangle = |\bar{d}s\rangle.
\]
CP eigenstates:
\[
|K_{S}\rangle = \frac{1}{\sqrt{2}}\bigl(|K^{0}\rangle + |\bar{K}^{0}\rangle\bigr),\qquad
|K_{L}\rangle = \frac{1}{\sqrt{2}}\bigl(|K^{0}\rangle - |\bar{K}^{0}\rangle\bigr).
\]
Evolution with non-Hermitian Hamiltonian
\[
\hat{H} = \mat{E_{S}-i\Gamma_{S}/2 & 0 \\ 0 & E_{L}-i\Gamma_{L}/2}
\]
gives the oscillatory transition probability
\[
\bigl|\langle\bar{K}^{0}\mid\Psi_{K}(t)\rangle\bigr|^{2}
= \frac{1}{4}\Bigl(e^{-t/\tau_{S}} + e^{-t/\tau_{L}}
- 2e^{-t/\tau^{*}}\cos\bigl((\omega_{S}-\omega_{L})t\bigr)\Bigr).
\]

% ============================================================
\section{Cabibbo Angle and See-Saw Mechanism}
\label{sec:1.4}

Flavour mixing in the $s$--$d$ sector:
\[
\mat{d' \\ s'} = e^{\vartheta\hat{i}}\mat{d \\ s}
= \mat{\cos\vartheta & \sin\vartheta \\ -\sin\vartheta & \cos\vartheta}\mat{d \\ s}.
\]
See-saw neutrino mass matrix:
\[
\hat{m}_{\nu} = \mat{0 & m_{D} \\ m_{D} & M_{NHL}},\qquad M_{NHL}\gg m_{D}.
\]

% ============================================================
\section{Gell-Mann and Pauli Matrices}
\label{sec:1.5}

\subsection{Spin Composition}

\[
\tfrac{1}{2}\otimes\tfrac{1}{2} = 0\oplus 1,\qquad
\tfrac{1}{2}\otimes\tfrac{1}{2}\otimes\tfrac{1}{2}
= \tfrac{1}{2}\oplus\tfrac{1}{2}\oplus\tfrac{3}{2}.
\]

\subsection{Gell-Mann Matrices}

The eight $3\times 3$ Gell-Mann matrices $\hat{\lambda}_{\alpha}$ ($\alpha=1,\dots,8$) satisfy
\[
[\hat{\lambda}_{\alpha},\hat{\lambda}_{\beta}] = i\,f_{\alpha\beta\gamma}\,\hat{\lambda}_{\gamma},\qquad
\tr(\hat{\lambda}_{\alpha}\hat{\lambda}_{\beta}) = 2\delta_{\alpha\beta}.
\]
Explicitly,
\begin{align*}
\hat{\lambda}_{1}&=\mat{0&1&0\\1&0&0\\0&0&0},&
\hat{\lambda}_{2}&=\mat{0&-i&0\\i&0&0\\0&0&0},&
\hat{\lambda}_{3}&=\mat{1&0&0\\0&-1&0\\0&0&0},\\
\hat{\lambda}_{4}&=\mat{0&0&1\\0&0&0\\1&0&0},&
\hat{\lambda}_{5}&=\mat{0&0&-i\\0&0&0\\i&0&0},&
\hat{\lambda}_{6}&=\mat{0&0&0\\0&0&1\\0&1&0},\\
\hat{\lambda}_{7}&=\mat{0&0&0\\0&0&-i\\0&i&0},&
\hat{\lambda}_{8}&=\frac{1}{\sqrt{3}}\mat{1&0&0\\0&1&0\\0&0&-2}.
\end{align*}

\subsection{Flavours and Spin}

Mesons: $3\otimes\bar{3} = 1\oplus 8$.\\
Baryons: $3\otimes 3\otimes 3 = 1\oplus 8\oplus 8\oplus 10$.

Proton wave function (flavour $\times$ spin):
\[
|p\rangle = \frac{1}{2\sqrt{3}}\,
(duu\;\;udu\;\;uud)\,
\mat{2&-1&-1\\-1&2&-1\\-1&-1&2}\,
\mat{\downarrow\uparrow\uparrow \\ \uparrow\downarrow\uparrow \\ \uparrow\uparrow\downarrow}.
\]

\subsection{Colour and QCD}

Colour triplet $R,G,B$ makes baryons antisymmetric in colour:
\[
|\Omega^{-}\rangle = \frac{1}{\sqrt{6}}\sum_{ijk}\varepsilon_{ijk}\,|s_{i}s_{j}s_{k}\rangle.
\]
Gluons belong to the $SU(3)$ colour octet.

% ============================================================
\section{Concluding Remarks}
\label{sec:1.6}

\subsection{Vector Differential Equations and Matrices}
\label{sec:1.6.1}

Heisenberg-like equation for the Pauli vector:
\[
\frac{d}{dt}\vec{\sigma} = \vec{\Omega}\times\vec{\sigma}.
\]
Solution (Rodrigues rotation):
\begin{align}
\vec{\sigma}(t)
&= \cos(\Omega t)\,\vec{\sigma}_{0}
 + \sin(\Omega t)\,(\vec{n}\times\vec{\sigma}_{0}) \notag \\
&\quad + \bigl[1-\cos(\Omega t)\bigr]\,(\vec{n}\!\cdot\!\vec{\sigma}_{0})\,\vec{n}.
\label{eq:rodrigues}
\end{align}
Non-homogeneous case:
\[
\frac{d}{dt}\vec{\sigma} = \vec{\Omega}\times\vec{\sigma} + \vec{\gamma}
\;\Longrightarrow\;
\vec{\sigma}(t) = e^{t[\vec{\Omega}]}\vec{\sigma}_{0}
+ \int_{0}^{t} e^{(t-t')[\vec{\Omega}]}\,\vec{\gamma}\,dt'.
\]

\subsection{Matrices, Vector Equations and Rotations}
\label{sec:1.6.2}

For $\vec{S} = (S_{1},S_{2},S_{3})^{T}$,
\[
\frac{d}{dt}\mat{S_{1}\\S_{2}\\S_{3}}
= \mat{0&-\Omega_{3}&\Omega_{2}\\
        \Omega_{3}&0&-\Omega_{1}\\
       -\Omega_{2}&\Omega_{1}&0}
  \mat{S_{1}\\S_{2}\\S_{3}}.
\]
The evolution operator $\hat{U}(t)$ is the Rodrigues rotation matrix
\[
\hat{U}(t) = \hat{1} + \frac{\sin\Omega t}{\Omega}\,\hat{\Omega}
+ \frac{1-\cos\Omega t}{\Omega^{2}}\,\hat{\Omega}^{2},
\]
with $\Omega = \sqrt{\Omega_{1}^{2}+\Omega_{2}^{2}+\Omega_{3}^{2}}$.

\subsection{Cabibbo--Kobayashi--Maskawa Matrix}
\label{sec:1.6.3}

Generalisation to three families:
\[
\mat{d'\\s'\\b'} = e^{\hat{\Omega}}\mat{d\\s\\b},
\]
with $\Omega_{3}=\vartheta_{c}$, $\Omega_{2}=x\vartheta_{c}^{2}$, $\Omega_{1}=y\vartheta_{c}^{3}$
and $\vartheta_{c}\simeq 0.22$ rad.

\subsection{Frenet--Serret Equations}
\label{sec:1.6.4}

\[
\frac{d}{ds}\mat{T\\N\\B}
= \mat{0&\kappa&0\\-\kappa&0&\tau\\0&-\tau&0}
  \mat{T\\N\\B}.
\]

\subsection{Matrix, Rotations and Euler Angles}
\label{sec:1.6.5}

Euler-angle factorisation:
\begin{align*}
\hat{R}
&= e^{\alpha\hat{J}_{z}}\,e^{\beta\hat{J}_{y}}\,e^{\gamma\hat{J}_{z}} \\
&= \mat{\cos\alpha&-\sin\alpha&0\\\sin\alpha&\cos\alpha&0\\0&0&1}
   \mat{1&0&0\\0&\cos\beta&-\sin\beta\\0&\sin\beta&\cos\beta}
   \mat{\cos\gamma&-\sin\gamma&0\\\sin\gamma&\cos\gamma&0\\0&0&1}.
\end{align*}

\subsection{4-Vectors and Lorentz Transformations}
\label{sec:1.6.6}

Minkowski metric:
\[
\eta_{\mu\nu} = \mat{-1&0&0&0\\0&1&0&0\\0&0&1&0\\0&0&0&1}.
\]
Lorentz boost along $x$:
\[
\mat{ct'\\x'\\y'\\z'}
= \mat{\gamma&-\gamma\beta&0&0\\-\gamma\beta&\gamma&0&0\\0&0&1&0\\0&0&0&1}
  \mat{ct\\x\\y\\z},
\qquad \gamma = \frac{1}{\sqrt{1-\beta^{2}}}.
\]

\subsection{Dirac Matrices}
\label{sec:1.6.7}

In block form (using $2\times 2$ blocks $\hat{1},\hat{i},\hat{h},\hat{t}$):
\[
\gamma^{0} = \mat{\hat{1}&\hat{0}\\\hat{0}&-\hat{1}},\quad
\gamma^{1} = \mat{\hat{0}&\hat{h}\\-\hat{h}&\hat{0}},\quad
\gamma^{2} = -i\mat{\hat{0}&\hat{i}\\-\hat{i}&\hat{0}},\quad
\gamma^{3} = \mat{\hat{0}&\hat{t}\\-\hat{t}&\hat{0}}.
\]
And
\[
\gamma^{5} = i\gamma^{0}\gamma^{1}\gamma^{2}\gamma^{3}
= \mat{\hat{0}&\hat{1}\\\hat{1}&\hat{0}}.
\]

The Dirac equation:
\[
i\hbar\,\partial_{t}\Psi = \vec{\alpha}\!\cdot\!\vec{p}\,\Psi + \beta\,mc^{2}\Psi.
\]

Foldy--Wouthuysen transformation:
\[
\hat{U} = e^{\beta\vec{\alpha}\!\cdot\!\vec{p}\,\vartheta}
= \cos\theta + \beta\,\vec{\alpha}\!\cdot\!\vec{n}_{p}\sin\theta,\qquad
\vec{n}_{p} = \frac{\vec{p}}{p},
\]
with $\tan(2\vartheta) = p/(mc)$.

% ============================================================
\section*{Bibliography for Chapter 1}
\addcontentsline{toc}{section}{Bibliography for Chapter 1}

\begin{thebibliography}{99}

\bibitem{gantmaker} F.~R.\ Gantmaker, \emph{The Theory of Matrices}, Chelsea, New York, 1959.

\bibitem{lancaster} P.~Lancaster, M.~Tismenetsky, \emph{The Theory of Matrices with Applications}, Academic Press, 1985.

\bibitem{bronson1} R.~Bronson, \emph{Matrix Methods: An Introduction}, Academic Press, 1970.

\bibitem{bronson2} R.~Bronson, \emph{Schaum's Outline of Theory and Problems of Matrix Operations}, McGraw-Hill, 1989.

\bibitem{brown} W.~Brown, \emph{Matrices and Vector Spaces}, Marcel Dekker, 1991.

\bibitem{leonhardt} U.~Leonhardt, \emph{Essential Quantum Optics}, Cambridge University Press, 2000.

\bibitem{gasiorowicz} S.~Gasiorowicz, \emph{Quantum Physics}, Wiley, 1974.

\bibitem{goldstein} H.~Goldstein, \emph{Classical Mechanics}, 2nd ed., Addison-Wesley, 1980.

\bibitem{dattoli1990} G.~Dattoli, M.~Richetta, G.~Schettini, A.~Torre,
``Lie algebraic methods and solutions of linear partial differential equations'', \emph{J.\ Math.\ Phys.}\ \textbf{31}, 2856 (1990).

\bibitem{wei1963} J.~Wei, E.~Norman,
``Lie algebraic solution of linear differential equations'', \emph{J.\ Math.\ Phys.}\ \textbf{4}, 575 (1963).

\bibitem{vanderlugt} A.~Vander Lugt,
``Operational Notation for the Analysis and Synthesis of Optical Data Processing Systems'', \emph{Proc.\ IEEE} \textbf{54}, 1055 (1966).

\bibitem{stoler} D.~Stoler, ``Operator Methods in Physical Optics'', \emph{J.\ Opt.\ Soc.\ Am.}\ \textbf{71}, 334 (1981).

\bibitem{bacry} H.~Bacry, M.~Cadilhac, ``Metaplectic Group and Fourier Optics'', \emph{Phys.\ Rev.\ A} \textbf{23}, 2533 (1981).

\bibitem{sanchez} J.~Sanchez-Mandragon, K.~B.~Wolf, \emph{Lie Methods in Optics}, Springer, 1986.

\bibitem{dattoli1988} G.~Dattoli, A.~Torre, J.~C.~Gallardo,
``Operatorial methods in Optics'', \emph{Riv.\ Nuovo Cimento} \textbf{11}, 1 (1988).

\bibitem{torre2005} A.~Torre, \emph{Linear Ray and Wave Optics in Phase Space}, Elsevier, 2005.

\bibitem{feynman1957} R.~P.~Feynman, L.~P.~Vernon, R.~W.~Hellwarth,
``Geometrical Representation of the Schrodinger Equation to Solve Maser Problems'', \emph{J.\ Appl.\ Phys.}\ \textbf{28}, 49 (1957).

\bibitem{wu1957} C.~S.~Wu \emph{et al.}, ``Experimental Test of Parity Conservation in Beta Decay'', \emph{Phys.\ Rev.}\ \textbf{105}, 1413 (1957).

\bibitem{yariv} A.~Yariv, \emph{Quantum Electronics}, Wiley, 1985.

\bibitem{hallen} L.~Hallen, J.~H.~Eberly, \emph{Optical Resonance and Two-Level Atoms}, Wiley, 1975.

\bibitem{baym} G.~Baym, \emph{Lectures on Quantum Mechanics}, Benjamin, 1969.

\bibitem{cabibbo} N.~Cabibbo, ``Unitary Symmetry and Leptonic Decays'', \emph{Phys.\ Rev.\ Lett.}\ \textbf{10}, 531 (1963).

\bibitem{kobayashi} M.~Kobayashi, T.~Maskawa,
``CP Violation in the Renormalizable Theory of Weak Interaction'', \emph{Prog.\ Theor.\ Phys.}\ \textbf{49}, 652 (1973).

\bibitem{wolfenstein} L.~Wolfenstein, ``Parametrization of the Kobayashi-Maskawa Matrix'', \emph{Phys.\ Rev.\ Lett.}\ \textbf{51}, 1945 (1983).

\bibitem{dattoli1994} G.~Dattoli, ``Weak mixing angles, CP-violating phase and exponential parametrization of the Cabibbo-Kobayashi-Maskawa matrix'', \emph{Nuovo Cim.\ A} \textbf{107}, 1243 (1994).

\bibitem{mohapatra} R.~N.~Mohapatra, ``Physics of Neutrino Masses'', SLAC Summer Institute (SSI04), 2004.

\bibitem{gellmann} M.~Gell-Mann, \emph{The Eightfold Way}, Benjamin, 1964.

\bibitem{lichtenberg} D.~B.~Lichtenberg, \emph{Unitary Symmetry and Elementary Particles}, Academic Press, 1978.

\bibitem{peskin} M.~E.~Peskin, D.~V.~Schroeder, \emph{An Introduction to Quantum Field Theory}, Addison Wesley, 1995.

\bibitem{mangano} M.~L.~Mangano, ``Introduction to QCD'', in \emph{Proc.\ 1998 European School of High-Energy Physics}.

\bibitem{babusci2011} D.~Babusci, G.~Dattoli, E.~Sabia,
``Operational Methods and Lorentz Type Equations'', \emph{J.\ Phys.\ Math.}\ \textbf{3}, P110601 (2011).

\bibitem{varshalovich} D.~A.~Varshalovich, A.~N.~Moskalev, V.~K.~Khersonskii,
\emph{Quantum Theory of Angular Momentum}, World Scientific, 1988.

\bibitem{morin} D.~Morin, \emph{Relativity (Kinematics)}, Ch.~11, Harvard University, 2007.

\bibitem{good} R.~H.~Good Jr., ``Properties of Dirac Matrices'', \emph{Rev.\ Mod.\ Phys.}\ \textbf{27}, 187 (1955).

\bibitem{bjorken} S.~D.~Bjorken, S.~D.~Drell, \emph{Relativistic Quantum Mechanics}, McGraw-Hill, 1964.

\end{thebibliography}

\section*{Transcribed Notes (raw OCR, verbatim)}
\begin{verbatim}
M
ATHEMATICAL
METHODS
FOR PHYSICISTS
11315_9789811201578_TP.indd   1 18/9/19   11:14 AM
This page intentionally left blankThis page intentionally left blank
NEW JERSEY   •  LONDON   •  SINGAPORE   •  BEIJING   •  SHANGHAI   •  HONG KONG   •  TAIPEI   •  CHENNAI   •  TOKYO
World Scientific
Danilo Babusci
INFN, Italy
Giuseppe Dattoli
ENEA, Italy
Silvia Licciardi
ENEA, Italy
Elio Sabia
ENEA, Italy
M
ATHEMATICAL
METHODS
FOR PHYSICISTS
11315_9789811201578_TP.indd   2 18/9/19   11:14 AM
Published by
World Scientific Publishing Co. Pte. Ltd.
5 Toh Tuck Link, Singapore 596224
USA office:  27 Warren Street, Suite 401-402, Hackensack, NJ 07601
UK office:  57 Shelton Street, Covent Garden, London WC2H 9HE
British Library Cataloguing-in-Publication Data
A catalogue record for this book is available from the British Library.
MATHEMATICAL
METHODS
FOR
PHYSICISTS
Copyright © 2020 by World Scientific Publishing Co. Pte. Ltd.
All rights reserved. This book, or parts thereof, may not be reproduced in any form or by any means, electronic or
mechanical, including photocopying, recording or any information storage and retrieval system now known or to
be invented, without written permission from the publisher.
For photocopying of material in this volume, please pay a copying fee through the Copyright Clearance Center,
Inc., 222 Rosewood Drive, Danvers, MA 01923, USA. In this case permission to photocopy is not required from
the publisher.
ISBN
978-981-120-157-8
ISBN
978-981-120-268-1 (pbk)
For any available supplementary material, please visit
https://www
.worldscientific.com/worldscibooks/10.1142/11315#t=suppl
Desk Editor: Nur Syarfeena Binte Mohd Fauzi
Typeset by Stallion Press
Email: enquiries@stallionpress.com
Printed in Singapore
Syarfeena - 11315 - Mathematical Methods for Physicists.indd   1 18-09-19   11:44:13 AM
v
“If you intend to mount heavy mathematical artillery again during your com-
ing year in Europe, I would ask you not only not to come to Leiden, but if
possible not even to Holland, and just because I am really so fond of you and
want to keep it that way. But if, on the contrary, you want to spend at least
your first few months patiently, comfortably, and joyfully in discussions that
keep coming back to the same few points, chatting about a few basic questions
with me and our young people - and without thinking much about publishing -
then I welcome you with open arms!! ”
Eherenfest in a letter to Oppenheimer, Summer 1928
This page intentionally left blankThis page intentionally left blank
Contents
Abstract xiii
1 Matrices, Exponential Operators and Physical Applications 1
1.1 Introduction . . . . . . . . . . . . . . . . . . . . . . . . . . . . 1
1.2 Pauli Matrices . . . . . . . . . . . . . . . . . . . . . . . . . . . 6
1.3 Applications of 2 × 2 Matrices . . . . . . . . . . . . . . . . . . 12
1.3.1 Classical Optics: Ray Beam Propagation and ABCD
Law . . . . . . . . . . . . . . . . . . . . . . . . . . . . 12
1.3.2 Quantum Mechanics . . . . . . . . . . . . . . . . . . . 19
1.3.3 Particle Physics: Kaon Mixing . . . . . . . . . . . . . . 22
1.4 Cabibbo Angle and See-Saw Mechanism . . . . . . . . . . . . 29
1.5 Gell-Mann and Pauli Matrices . . . . . . . . . . . . . . . . . . 33
1.5.1 Spin Composition . . . . . . . . . . . . . . . . . . . . . 33
1.5.2 Gell-Mann Matrices . . . . . . . . . . . . . . . . . . . . 36
1.5.3 Flavors and Spin . . . . . . . . . . . . . . . . . . . . . 38
1.5.4 Color and QCD . . . . . . . . . . . . . . . . . . . . . . 41
1.6 Concluding Remarks . . . . . . . . . . . . . . . . . . . . . . . 43
1.6.1 Vector Differential Equations and Matrices . . . . . . . 44
1.6.2 Matrices, Vector Equations and Rotations . . . . . . . 47
1.6.3 Cabibbo-Kobayashi-Maskawa Matrix . . . . . . . . . . 50
1.6.4 Frenet-Serret Equations . . . . . . . . . . . . . . . . . 51
1.6.5 Matrix, Rotations and Euler Angles . . . . . . . . . . . 52
1.6.6 4-Vectors and Lorentz Transformations . . . . . . . . . 53
1.6.7 Dirac Matrices . . . . . . . . . . . . . . . . . . . . . . 55
Bibliography . . . . . . . . . . . . . . . . . . . . . . . . . . . . . . 59
vii
viii Mathematical Methods for Physicists
2 Ordinary and Partial Differential Equations, Evolution
Operator Method and Applications 63
2.1 Ordinary Differential Equations, Matrices and Exponential
Operators . . . . . . . . . . . . . . . . . . . . . . . . . . . . . 63
2.2 Partial Differential Equations and Exponential
Operators, I . . . . . . . . . . . . . . . . . . . . . . . . . . . . 66
2.3 Partial Differential Equations and Exponential
Operators, II . . . . . . . . . . . . . . . . . . . . . . . . . . . 72
2.4 Operator Ordering . . . . . . . . . . . . . . . . . . . . . . . . 73
2.5 Schr¨ odinger Equation and Paraxial Wave Equation
of Classical Optics . . . . . . . . . . . . . . . . . . . . . . . . 79
2.6 Examples of Fokker-Planck, Schr¨ odinger and Liouville
Equations . . . . . . . . . . . . . . . . . . . . . . . . . . . . . 84
2.7 Concluding Remarks . . . . . . . . . . . . . . . . . . . . . . . 88
Bibliography . . . . . . . . . . . . . . . . . . . . . . . . . . . . . . 95
3 Hermite Polynomials and Applications 101
3.1 Introduction . . . . . . . . . . . . . . . . . . . . . . . . . . . . 101
3.2 Hermite Polynomials Generating Function . . . . . . . . . . . 104
3.2.1 Introducing the Generating Function . . . . . . . . . . 104
3.2.2 Generating Function Applications . . . . . . . . . . . . 106
3.3 Hermite Polynomials as an Orthogonal Basis . . . . . . . . . . 109
3.4 Hermite Polynomials in Quantum Mechanics: Creation and
Annihilation Operators . . . . . . . . . . . . . . . . . . . . . . 113
3.5 Quantum Mechanics Applications . . . . . . . . . . . . . . . . 117
3.6 Coherent or Quasi-Classical States of Harmonic Oscillators . . 121
3.7 Jaynes-Cummings Model . . . . . . . . . . . . . . . . . . . . . 127
3.8 Classical Optics and Hermite Polynomials . . . . . . . . . . . 129
Bibliography . . . . . . . . . . . . . . . . . . . . . . . . . . . . . . 132
4 Laguerre Polynomials, Integral Operators and
Applications 135
4.1 Introduction . . . . . . . . . . . . . . . . . . . . . . . . . . . . 135
4.2 Laguerre Polynomials Generating Function . . . . . . . . . . . 140
4.3 Orthogonality Properties of Laguerre Polynomials . . . . . . . 141
4.4 Bessel Functions . . . . . . . . . . . . . . . . . . . . . . . . . 144
4.5 Associated Laguerre Polynomials . . . . . . . . . . . . . . . . 148
4.6 Legendre Polynomials . . . . . . . . . . . . . . . . . . . . . . . 151
Contents ix
4.7 Miscellaneous Applications and Comments . . . . . . . . . . . 153
4.8 App´ el Polynomials and Final Comments . . . . . . . . . . . . 158
Bibliography . . . . . . . . . . . . . . . . . . . . . . . . . . . . . . 163
5 Exercises and Complements I 167
5.1 Pauli and Jones Matrices and Mueller Calculus . . . . . . . . 167
5.2 Magnetic Lenses and Matrix Description . . . . . . . . . . . . 174
5.3 Miscellanea on the Matrix Formalism and Solution
of Evolution Problems . . . . . . . . . . . . . . . . . . . . . . 182
5.3.1 Matrices and Quaternions . . . . . . . . . . . . . . . . 182
5.3.2 Matrix Solution of Evolution Problems . . . . . . . . . 183
5.4 Lorentz Transformation . . . . . . . . . . . . . . . . . . . . . . 186
5.5 Hyperbolic Trigonometry and Special Relativity . . . . . . . . 189
5.6 A Touch on Elliptic Functions . . . . . . . . . . . . . . . . . . 194
5.7 Concluding Comments . . . . . . . . . . . . . . . . . . . . . . 207
Bibliography . . . . . . . . . . . . . . . . . . . . . . . . . . . . . . 212
6 Exercises and Complements II 217
6.1 Ordinary Differential Equations and Matrices . . . . . . . . . 217
6.2 Crofton-Glaisher Identities and Heat Type Equations . . . . . 230
6.3 Gamma Function and Definite Integrals . . . . . . . . . . . . . 235
6.4 Complex Variable Method and Evaluation of Integrals . . . . . 244
6.5 Fourier Transform . . . . . . . . . . . . . . . . . . . . . . . . . 250
6.6 Fourier Transform and the Solution of Differential
Equations . . . . . . . . . . . . . . . . . . . . . . . . . . . . . 259
6.7 Fourier-Type Transforms . . . . . . . . . . . . . . . . . . . . . 262
Bibliography . . . . . . . . . . . . . . . . . . . . . . . . . . . . . . 267
7 Exercises and Complements III 273
7.1 Second Solution of Hermite Equation . . . . . . . . . . . . . . 273
7.2 Higher Orders Hermite Polynomials . . . . . . . . . . . . . . . 275
7.3 Multi-Index Hermite Polynomials . . . . . . . . . . . . . . . . 280
7.4 Creation-Annihilation Operators Algebra and Physical
Applications . . . . . . . . . . . . . . . . . . . . . . . . . . . . 284
7.5 Eisenstein Integers . . . . . . . . . . . . . . . . . . . . . . . . 290
7.6 Harmonic Oscillator Hamiltonian Formal Aspects and Further
Miscellaneous Considerations . . . . . . . . . . . . . . . . . . 295
7.7 Time-Dependent Hamiltonians . . . . . . . . . . . . . . . . . . 298
x Mathematical Methods for Physicists
7.8 Dyson Series . . . . . . . . . . . . . . . . . . . . . . . . . . . . 299
7.8.1 Beyond the Dyson Expansion . . . . . . . . . . . . . . 303
7.9 Special Polynomials and Perturbation Theory . . . . . . . . . 307
Bibliography . . . . . . . . . . . . . . . . . . . . . . . . . . . . . . 311
8 Exercises and Complements IV 319
8.1 Sturm-Liouville Problem . . . . . . . . . . . . . . . . . . . . . 319
8.2 Green’s Functions . . . . . . . . . . . . . . . . . . . . . . . . . 324
8.3 Laguerre Polynomials, Associated Operators and PDE . . . . 326
8.4 App´ el Polynomials, Associated Operators and Partial
Differential Equations . . . . . . . . . . . . . . . . . . . . . . . 331
8.5 Riemann Function . . . . . . . . . . . . . . . . . . . . . . . . 337
8.6 Bessel Special Functions . . . . . . . . . . . . . . . . . . . . . 342
Bibliography . . . . . . . . . . . . . . . . . . . . . . . . . . . . . . 351
9 Special Functions, Umbral Methods and Applications 355
9.1 Introduction to Umbral Methods and Relevant
Applications . . . . . . . . . . . . . . . . . . . . . . . . . . . . 355
9.2 Further Comments on Umbral Methods, Infinite Integrals
and Borel Transform . . . . . . . . . . . . . . . . . . . . . . . 362
9.3 Borel Transform and Applications . . . . . . . . . . . . . . . . 368
9.3.1 Borel - Pad´ e . . . . . . . . . . . . . . . . . . . . . . . . 371
9.4 Umbral Formalism and Laguerre Polynomials . . . . . . . . . 378
9.5 Umbral Formalism and Hermite Polynomials . . . . . . . . . . 380
9.6 Umbral Formalism and Operator Ordering . . . . . . . . . . . 382
9.7 Mittag-Leffler Function and Fractional Calculus
Application . . . . . . . . . . . . . . . . . . . . . . . . . . . . 386
9.8 Formalism of Negative Derivative and Definite Integrals . . . . 392
9.9 Umbral Formalism, Dual Numbers and Super-Gaussian
Beam Transport . . . . . . . . . . . . . . . . . . . . . . . . . . 394
Bibliography . . . . . . . . . . . . . . . . . . . . . . . . . . . . . . 403
10 A Glimpse into the Math of the Feynman Diagrams 411
10.1 Introduction Non Relativistic Scattering Theory
and Lippman-Schwinger Equation . . . . . . . . . . . . . . . . 411
10.2 Fermi Golden Rule . . . . . . . . . . . . . . . . . . . . . . . . 417
10.2.1 Fermi Golden Rule Application . . . . . . . . . . . . . 421
10.3 Feynman Diagrams: Introductory Rules . . . . . . . . . . . . . 423
Contents xi
10.4 Virtual Particles and Propagators . . . . . . . . . . . . . . . 427
10.5 Space and Time like Feynman Diagrams . . . . . . . . . . . . 430
10.6 Dirac Gamma Matrices . . . . . . . . . . . . . . . . . . . . . . 431
10.7 Mathematics of the Dirac Equation . . . . . . . . . . . . . . . 439
10.8 A Touch on Quantum Electrodynamics . . . . . . . . . . . . . 445
10.9 Formal Point of View to the Dimensions and Units
in Physics . . . . . . . . . . . . . . . . . . . . . . . . . . . . . 450
Bibliography . . . . . . . . . . . . . . . . . . . . . . . . . . . . . . 460
Index 463
This page intentionally left blankThis page intentionally left blank
Abstract
This book reports on Mathematical Methods for Physics elaborated
on the original plots of the lectures given by one of the Authors, Giuseppe
Dattoli, at the Rome “Tre” University in the academic year 2008-2009.
During these laps of time, the lectures have appeared in the form of an
internal note http://opac22.bologna.enea.it/RT/2010/2010 58 ENEA.pdf co-
authored by Danilo Babusci, Giuseppe Dattoli and Mario Del Franco. The
present version is completely revisited and new topics have been included. In
particular two new chapters have been added.
The book plane follows the orginal conception guided by the question:
“Which methods for which Physics?”
Our goal is not that of providing a series of lectures on well known problems
in Mathematical Physics but, rather, to give some tools to attack specific
problems which are often encountered in different branches of Physics, ran-
ging from Quantum Field Theory, to Laser Physics and to Charged Particle
Transport in Accelerators.
The power of Mathematics, or better what has been called its “unreaso-
nable success” in describing physical phenomena, is just the fact that the same
tools can be exploited to model different aspects of Physics. The formalism of
matrices is useful in Quantum Mechanics to study the dynamical behavior of
particles with spin but it turns out to be useful to study the propagation of
optical beams through lens systems or to treat the charged beam transport in
Linacs or in Storage Rings. The Schr¨ odinger equation can be used to study
the quantum mechanical spreading of a free particle and the distribution of
electromagnetic fields in an optical cavity as well. The heat equation can be
xiii
xiv Mathematical Methods for Physicists
exploited in so a large number of contexts which even includes the behavior
of the stock options in the financial markets.
Therefore, our choice has been to pick up a certain number of mathemat-
ical techniques, discuss them carefully and present the physical problems to
which they apply. We have touched on problems in Classical and Quantum
Electromagnetism, in particle and solid state Physics, by taking into account
a kind of unitarity (at least formal) between the different physical phenomena.
This is an attempt which requires some effort by the reader, since we
have not used sharp cuts between Quantum and Classical Mechanics or be-
tween high energy and solid state Physics, we have mainly been guided by the
evolution and by the internal mathematical consistency of ideas flowing from
matrices, to differential operators, to group theory . . .
We owe our gratitude to Prof. Orlando Ragnisco and to Prof. Decio Levi
who allowed this experiment and encouraged and advised us during the prepa-
ration of the original version of the book.
The elaboration of the present version comes after years of intensive work
on different topics in Theoretical Physics and Mathematics, including a new
point of view to Special Functions via Umbral Methods and research on elec-
tromagnetic processes ruled by non local equations.
We have therefore spread out these further experience in the different
chapters. We have, for example, dedicated chapter IX (written by Giuseppe
Dattoli and Silvia Licciardi) to the use of umbral methods, including ap-
plication to Physics. Moreover, we have added the chapter X (written by
Giuseppe Dattoli and Elio Sabia) to provide a less qualitiative understanding
of the mathematical aspects underlying the technicalities of Feynman Dia-
grams. This chapter has benefitted from the kind help of Dr. Federico Nguyen
who corrected the original draft.
It is finally a pleasure to recognize the interest of the late Prof. Benedetto
De Tollis who gave access to his unpublished lectues notes on Feynman Dia-
grams given at Rome “La Sapienza” University.
Chapter 1
Matrices, Exponential
Operators and Physical
Applications
1.1 Introduction
The matrices and the relevant formalism are widely exploited tools in
Physics and Mathematics. They provide the back bone of the treatment of
many problems in classical and quantum Physics. Thus, we start these lec-
tures with a short resume on the relevant properties, aimed at fixing the
notation and at illustrating the technicalities we will later employ in solving
mathematical problems of physical interest.
Our starting point is the equation 1
ˆAζ =b (1.1.1)
where
ˆA =
( α1 1 α1 2
α2 1 α2 2
)
, ζ =
( z1
z2
)
, b =
( b1
b2
)
, ∀αij, zi, bi∈ C.
(1.1.2)
Eq. (1.1.1) can be considered as the generalization of an algebraic equation of
first degree (it represents indeed an algebraic system of first degree equations)
1The symbol ˆO will denote either a matrix or an operator, the column vectors will be
denoted by δ , the scalars will not be specified by any extra symbol.
1
2 Mathematical Methods for Physicists
and, provided that we accept the matrix as the generalization of the ordinary
numbers, its solution can, therefore, be written in the form
ζ = ˆA−1b, (1.1.3)
reminiscent of the solution of first degree algebraic equations. Such a formal
solution makes sense if the matrix ˆA is not singular or, what is the same, if
it is invertible. It goes by itself that, in deriving eq. (1.1.3), we have tacitly
assumed that
ˆA ˆA−1 = ˆ1 (1.1.4)
where we have denoted with
ˆ1 =
( 1 0
0 1
)
(1.1.5)
the real-unity matrix, satisfying the obvious property
ˆ1 ˆB = ˆB ˆ1 = ˆB, ˆ1ζ =ζ, (1.1.6)
valid for any generic matrix. Regarding the two dimensional case, it is easily
checked that the inverse of ˆA writes
ˆA−1 = 1
∆
( α2 2 −α1 2
−α2 1 α1 1
)
(1.1.7)
where ∆ is the matrix determinant
∆ =α1 1α2 1−α1 2α21. (1.1.8)
A matrix will be said singular, and thus not invertible, if its determinant is
vanishing.
Needless to say, we have just “rediscovered” the well known solution of a
system of first order algebraic equation with the Cramer’s method by using
a point of view aimed at showing that, within certain limits, we can use
matrices as suitable generalizations of the ordinary numbers (complex or
real). We can therefore define simple procedures to introduce operations like
sum, difference, product and power. As a further support to the previous
statements, let us note that, along with the matrix (1.1.5), playing the role
of unity, we can introduce the matrix
ˆi =
( 0 1
−1 0
)
(1.1.9)
Matrices, Exponential Operators and Physical Applications 3
called the imaginary matrix (or sympletic unit ), which may be understood
as a kind of “imaginary” unity matrix, since
ˆi2 =−ˆ1. (1.1.10)
We can make a further step in the direction of considering matrices as
ordinary numbers, by introducing the notion of the function of a matrix .
Accordingly, we will say that
ˆF =f( ˆA) (1.1.11)
is a matrix, which is explicitly defined through the Mac Laurin series expan-
sion of the function f(x), namely
f( ˆA) =
∞∑
n=0
f (n)(0)
n!
ˆAn (1.1.12)
with f (n)(0) denoting the nth derivative calculated at x = 0. The above
expansion is just a formal expression, since we did not specify any radius
of convergence which, on the other side, cannot be stated from the expan-
sion itself. The convergence can be checked a posteriori using, for example,
the following procedure. By considering the matrix operator ˆF acting on a
generic column vector ζ, we have
ˆF ζ=
∞∑
n=0
f (n)(0)
n!
ˆAnζ =
∞∑
n=0
f (n)(0)
n! ζn, ζ n = ˆAnζ. (1.1.13)
It is evident that, once we have defied the vector ζn, we can use the ordinary
criteria to establish whether the series in (1.1.13) is convergent. Going back
to eq. (1.1.12) we can also state that a sufficient (but not necessary) condi-
tion to define the function of a matrix is that the function admits a series
expansion. In less na¨ ıve terms, we can also say that a function of a matrix
can be defined if the function is “locally” given by a convergent power series.
In other words, we require that f(x) be at least analytic. The adjective local
(reported between quotes) will be clarified later.
A simple but important example is offered by the “exponentiation” of
the imaginary unit matrix, which will (nicely) bring us to a generalization
4 Mathematical Methods for Physicists
of the Euler formula. According to the previous discussion, the exponential
function having ˆi as argument, can be handled as it follows
eαˆi =
∞∑
n=0
αn
n!
ˆin =
∞∑
n=0
(−1)nα2n
(2n)!
ˆ1 +
∞∑
n=0
(−1)nα2n+1
(2n + 1)!
ˆi = cosα ˆ1 + sinαˆi.
(1.1.14)
The use of the correspondences
1→ ˆ1, i → ˆi (1.1.15)
and of the cyclical identities
ˆin =
{ (−1)mˆ1 n = 2m,
(−1)mˆi n = 2m + 1, (1.1.16)
suggest that eq. (1.1.14) is a generalization of the Euler formula , involving
the imaginary unit matrix instead of the imaginary unit. Written in a matrix
form, eq. (1.1.14) yields
eαˆi = ˆR(α) =
( cosα sinα
− sinα cosα
)
. (1.1.17)
We can therefore state the following.
Theorem 1. The exponentiation of the imaginary matrix leads to the two
dimensional rotation matrix.
We have underscored that the matrices are generalizations of the ordi-
nary numbers and they have therefore more general properties which will be
discussed below.
Let us now consider the 2× 2 matrix
ˆh =
( 0 1
1 0
)
(1.1.18)
which will be defined the hyperbolic unit matrix. It is, indeed, easily checked
that it satisfies the relation
ˆh2 = ˆ1 (1.1.19)
Matrices, Exponential Operators and Physical Applications 5
and that its exponentiation yields
eα ˆh = ˆRh(α) =
( coshα sinhα
sinhα coshα
)
. (1.1.20)
Note that the matrix (1.1.20) represents a rotation in a hyperbolic space
(hence the name hyperbolic of the unit ˆh) and its mathematical and physi-
cal meaning will be more deeply discussed in the following. Unlike ordinary
numbers, matrices are not commuting entities. We may, thus, expect that
the three units “ ˆ1, ˆi, ˆh” do not commute. It is easily checked that the real
unit matrix commutes, by definition, with any matrix, while the imaginary
and hyperbolic units yield
[
ˆi, ˆh
]
= ˆi ˆh− ˆhˆi = 2 ˆt, ˆt =
( 1 0
0 −1
)
. (1.1.21)
The matrix ˆt is a fourth unit and satisfies the conditions
eα ˆt = ˆS(α) =
( eα 0
0 e−α
)
, ˆt2 = ˆ1. (1.1.22)
The matrix ˆS(α) is called the squeezing mapping. Its role in physical prob-
lems is of paramount importance and will be discussed later. The four ma-
trices ˆi, ˆh, ˆt, ˆ1 represent the basis to express any 2× 2 matrix,( α1 1 α12
α21 α22
)
=α ˆ1 +β ˆi +γ ˆh +δ ˆt, (1.1.23)
where α, β, γ, δare the coefficients of the linear combination and can be
expressed in terms of the matrix entries. If we limit our analysis to traceless
matrices only, i.e. it is zero the sum of the elements of the principal diagonal
α1 1 +α22 = 0, (1.1.24)
we can trace out the real unity matrix and use the remaining units as basis.
Before concluding this section, let us remind that we can define the her-
mitian conjugate of the matrix ˆA as
ˆA+ =
( α∗
11 α∗
21
α∗
12 α∗
22
)
(1.1.25)
and find that the simplectic unit is anti-hermitian
ˆi+ =−ˆi, (1.1.26)
while all the other matrices are Hermitian.
6 Mathematical Methods for Physicists
1.2 Pauli Matrices
In the previous section, we have introduced a set of matrices and we have
shown that at least two of them, once exponentiated, generate rotations of
circular or hyperbolic nature. Concerning the circular case, we have stressed
that the operation
ψ =eαˆiξ ⇒
( x′
y′
)
=
( cosα sinα
− sinα cosα
)( x
y
)
(1.2.1)
generates the rotation of the plane coordinates x,y . This conclusion, even
though apparently trivial, is extremely interesting, and opens the way to
further speculations about the notion of Lie Group, which will be discussed
later. By keeping the derivative of both sides with respect to α, we find
dψ
dα = ˆiψ. (1.2.2)
Furthermore if we set
ˆi = ˆσ+− ˆσ−, α = Ωτ (1.2.3)
with
ˆσ+ =
( 0 1
0 0
)
, ˆσ− =
( 0 0
1 0
)
(1.2.4)
satisfying the identity ˆσ2
± = ˆ0 (with ˆ0 being the null matrix, i.e. a matrix
with all vanishing entries), we can recast eq. (1.2.2) in the form
dψ
dτ = Ω (ˆσ+− ˆσ−) ψ. (1.2.5)
This is essentially the Schr¨ odinger equationwith the Hamiltonian
ˆH =i ℏ Ω (ˆσ+− ˆσ−). (1.2.6)
The “hermiticity” of ˆH is ensured by the fact that the operators (1.2.4)
satisfies also the identitiy
ˆσ+
± = ˆσ∓. (1.2.7)
As we will see in the forthcoming parts of this chapter, the Hamiltonian
(1.2.6) is a paradigmatic tool for the study of two-level quantum mechanical
systems. However, before discussing these aspects of the problem, some na¨ ıve
Matrices, Exponential Operators and Physical Applications 7
algebraic manipulations are in order. We first note that the above matrices
are not commuting, we get
[ˆσ+, ˆσ−] = 2 ˆσ3, ˆσ3 = 1
2
ˆt =
( 1
2 0
0 − 1
2
)
, [ ˆσ3, ˆσ±] =± ˆσ±
(1.2.8)
which are equivalent to the commutation relations of the angular momentum
operator components in Quantum Mechanics. By using the language of the
theory of groups, we can say that these matrices realize one of the repre-
sentations of the generators of the uni-modular group SU (2)2. The matrices
ˆσ±, ˆσ3 are the Pauli matrices3, introduced during the pioneering days of
Quantum Mechanics to describe the spin dynamics. We can now apply the
previous discussion to the solution of the Schr¨ odinger problem (1.2.5) which
can be written in the form
ψ(τ) = ˆU(τ)ψ0 (1.2.9)
where the operator
ˆU(τ) =eΩτ(ˆσ+−ˆσ−) (1.2.10)
denotes the evolution operator associated with the solution of eq. (1.2.5)
which, from the mathematical point of view, is an initial value Cauchy pro-
blem. The unitarity4 of the ˆU operator is easily checked and follows from
the hermiticity of the operator (1.2.6). The argument of the exponential in
(1.2.10) contains two matrices which are not commuting quantities. It is
therefore evident that this imposes some cautions in handling the so called
disentanglement of the above exponential as a product of exponential opera-
tors. We have indeed that, as an obvious consequence of non commutative
character of the matrices, the “semigroup” property 5 of the exponential
eaeb =ebea =ea+b (1.2.11)
2In the following chapter we will discuss different realizations of the SU (2) group in-
cluding the more familiar case employing differential operators.
3The Pauli matrices standard definition does not provide for ˆσ± but ˆσ1 = ˆh, ˆσ2 =−iˆi,
ˆσ3 = ˆt.
4We remind that an operator is said unitary if its hermitian conjugate is equal to its
inverse. From eq. (1.2.10) we get indeed ˆU+(τ) =eΩτ (ˆσ+
+−ˆσ+
−) =e−Ωτ (ˆσ+−ˆσ−) = ˆU−1(τ).
5In Mathematics, a semigroup is an algebraic structure consisting of a set closed under
a binary associative operation.
8 Mathematical Methods for Physicists
does not hold anymore, and
eΩτ(ˆσ+−ˆσ−)̸=eΩτ ˆσ+eΩτ ˆσ−̸=eΩτ ˆσ−eΩτ ˆσ+. (1.2.12)
This fact implies that we must take into account appropriate ordering crite-
ria. The problems underlying operator ordering play a central role in modern
Theoretical Physics and will carefully be treated in this book. Here we use
a simple but effective argument which can be easily generalized. We assume
that a convenient ordering for the previous exponential operator is a product
of three exponential operators, containing the group generators, namely
eα (ˆσ+−ˆσ−) =e2h(α) ˆσ3ef(α) ˆσ+e−g(α) ˆσ−, α = Ωτ, (1.2.13)
where the functions h(α), f(α), g(α) are the characteristic ordering func-
tions, satisfying the “initial” condition
h(0) =f(0) =g(0) = 0. (1.2.14)
To proceed further we note that we have, as a consequence of the condition
U(0) = 1,
ef(α) ˆσ+ =
( 1 f(α)
0 1
)
, e −g(α) ˆσ− =
( 1 0
−g(α) 1
)
,
e2h(α)ˆσ3 =
( eh(α) 0
0 e−h(α)
) (1.2.15)
and therefore the products of the previous matrices yields
e2h(α) ˆσ3ef(α) ˆσ+e−g(α) ˆσ− =
( eh(α) 0
0 e−h(α)
)( 1 f(α)
0 1
)( 1 0
−g(α) 1
)
=
( eh(α)(1−g(α)f(α)) f(α)eh(α)
−g(α)e−h(α) e−h(α)
)
.
(1.2.16)
It is useful to verify for Exercise that the above matrix has unit determi-
nant. Its unitarity is a consequence of the fact that the operator α (ˆσ+− ˆσ−)
is anti-Hermitian6. Furthermore, because
Ω τ (ˆσ+− ˆσ−) = Ωτ ˆi, (1.2.17)
6It is possible also prove that the unitarity condition imposes that 1
f(α)− 1
g(α) =g(α).
Matrices, Exponential Operators and Physical Applications 9
we know that the overall effect of the evolution operator is that of produc-
ing a rotation, which can always be decomposed in three successive distinct
actions, generated by the exponential operators. We will discuss the physi-
cal meaning of the previous result in the forthcoming sections. Here we note
that by comparing eq. (1.2.16) with the rotation matrix, we can specify the
characteristic functions as follows 7
eh(α) = secα, f (α) = 1
2 sin 2α, g (α) = tanα. (1.2.18)
In Fig. 1.1 we have reported a geometric interpretation of the characteristic
functions,which are related to a kind of decomposition of the ordinary rota-
tion matrix. Further comments on such interpretation can be found in the
second part of this book (sec 5.5).
Figure 1.1: Characteristic functions g(α),f (α) and eh(α).
The previous results are relevant to the particular case in which the matrix
in the argument of the exponential contains non diagonal entries only. We
will treat a more general situation in which the matrix is anti-hermitian
and with vanishing trace. To this aim, we introduce the following operator
written in vector form ⃗ σ= (ˆσx, ˆσy, ˆσz) with
ˆσx = 1
2(ˆσ+ + ˆσ−) = 1
2
( 0 1
1 0
)
, ˆσy = 1
2i(ˆσ+− ˆσ−) = 1
2i
( 0 1
−1 0
)
,
ˆσz = ˆσ3 = 1
2
( 1 0
0 −1
)
.
(1.2.19)
7Note that∀A,B matrices we have A =B if a11 =b11,a 21 =b21, ... .
10 Mathematical Methods for Physicists
The vector components satisfy the relations of commutation of angular mo-
mentum
[ˆσj, ˆσk] =iεj,k,l ˆσl, j, k, l ≡x, y,z, (1.2.20)
where εj,k,l is the Levi Civita tensor . The use of such an operator allows
the definition of the following Hamiltonian operator
ˆH = ℏ⃗Ω·⃗ σ= ℏ Ω⃗ n·⃗ σ, ⃗ n =
⃗Ω
Ω (1.2.21)
where Ω is the modulus of the vector ⃗Ω, usually called the torque vector,
and ⃗ nis the unit vector defining its direction. The physical meaning of the
previous Hamiltonian will be discussed in the forthcoming sections.
The associated evolution operator ˆU =e−i ˆH
ℏ t writes
ˆU(t) =e−i Ωt⃗ n·⃗ σ (1.2.22)
and the reader will not find any problem in proving that
ˆU(t) = cos
(
Ω t
2
)
ˆ1− 2i sin
(
Ωt
2
)
⃗ σ·⃗ n (1.2.23)
(hint: note that ( ⃗ σ·⃗ n)2 = 1
4
ˆ1 and use the same procedure leading to eq.
(1.1.14)).
Before concluding this section we want to stress the following points:
a) An exponential operator containing a 2 × 2 matrix (with vanishing
trace) can always be written as the ordered product of three exponen-
tials containing the Pauli matrices, like in eq. (1.2.13) 8
ˆU(t) =eiωτ ˆσ3+Ωt (ˆσ+−ˆσ−) =e2h(α) ˆσ3ef(α) ˆσ+e−g(α) ˆσ− . (1.2.24)
8The evolution operator is associated with a Hamiltonian operator
ˆH =−ℏω ˆσ3 +i ℏ Ω (ˆσ+− ˆσ−).
Matrices, Exponential Operators and Physical Applications 11
b) The ordering (1.2.13) is valid for the exponentiation of any linear com-
bination of the angular momentum operators and is independent of the
specific representation (this aspect of the problem will be discussed in
the forthcoming Chapters).
It can be useful to prove the degree of understanding of the topics dis-
cussed in these sections by considering the problems proposed below.
Exercise 1. i) The operator ordering in eq. (1.2.13) is one of the 3 possible
ordered forms, it goes by itself that each form will be characterized by an
appropriate set of characteristic functions as e.g.
ˆU(τ) =e−γ(α) ˆσ−eφ(α) ˆσ+e2η(α) ˆσ3. (1.2.25)
Find the link between the different characteristic functions, corresponding to
different exponential ordering.
ii) Proof the identity
⃗ σ×⃗ σ=i⃗ σ (1.2.26)
and the understanding of its physical and geometric meaning
(hint: use the ordinary definition of the vector product and then the commu-
tation relations (1.2.8)).
iii) Finally, given a generic 2× 2 matrix ˆα with trace̸= 0, prove that
exp
(( α11 α12
α21 α2 2
))
=
( A11 A12
A2 1 A22
)
, (1.2.27)
where
A1 1 =
{
1√
∆
(α1 1−α2 2) sinh
(√
∆
2
)
+ cosh
(√
∆
2
)}
e
1
2Tr (ˆα),
A2 2 =
{
− 1√
∆
(α1 1−α2 2) sinh
(√
∆
2
)
+ cosh
(√
∆
2
)}
e
1
2Tr (ˆα),
A12
α1 2
= A2 1
α2 1
=
{
2√
∆
sinh
(√
∆
2
)}
e
1
2Tr (ˆα),
Tr(ˆα) =α11 +α2 2, ∆ = (α11−α22)2 + 4α1 2α2 1
(1.2.28)
12 Mathematical Methods for Physicists
(hint: take into account that the matrix at the exponent is not traceless there-
fore it writes as a combination of Pauli matrices including the unit matrix).
Here we have introduced and discussed the properties of 2 × 2 matrices
useful to describe two-level systems or spin doublets (or spin- 1
2 particles).
The case of three level systems or spin-1 particles (spin triplets) requires the
use of 3× 3 matrices, as it will be briefly discussed in the next sections and
more carefully in Chapter 5.
The results obtained in this section will provide the mathematical appa-
ratus for treating the physical problems we will discuss in the forthcoming
sections.
1.3 Applications of 2 × 2 Matrices
1.3.1 Classical Optics: Ray Beam Propagation and
ABCD Law
Talking about Pauli matrices, the most natural application that comes
to mind is the dynamics of particles with spin- 1
2 in a magnetic field. This
problem will be discussed in the following section. Here we show how the
Pauli matrices can be used to discuss other physical problems, usually not
treated within these frameworks.
Let us consider the transmission of optical beams through an etalon,
which can be considered as an interface of thickness l between media with
different optical indices. This process is usually described using elementary
mathematical tools, however the propagation of optical beams through such
an optical device can be conveniently described using the matrix formalism
which may provide, as we will see below, a significant improvement simplify-
ing the associated computational tools. The geometry of the ray transmission
and reflection inside the etalon is shown in Fig. 1.2. The direction of prop-
agation is specified by the dashed line perpendicular to the surfaces. The
ray position and angles (see Figs. 1.2-1.5) are defined with respect to the
optical axis and here we limit ourselves to the case of paraxial propagation,
i.e. ray propagation involving small angular deviation from the propagation
axis. If we characterize the ray beam in terms of a column vector having
Matrices, Exponential Operators and Physical Applications 13
Figure 1.2: Optical ray beams reflected and transmitted through an etalon.
its components as the distance from the optical axis and the angle it forms
with the axis direction, the problem of the ray propagation 9, whatever com-
plicated the optical system is, can always be reduced to that of finding the
relationship between input and output rays, in terms of the relation
( r
ϑ
)
0
=
( A B
C D
)( r
ϑ
)
i
(1.3.1)
where the matrix with entries A,B,C,D is the transport matrix and charac-
terizes the optical element (or the combination of a series of optical elements)
guiding the beam.
Figure 1.3: Input and output ray-beams from a complicated optical system. The dotted
line represents the propagation axis and position and angles are taken with respect to it
as indicated in the figure.
9Along with propagation we will use as “synonyms” evolution and transport.
14 Mathematical Methods for Physicists
The simplest optical transport element is the free space10 (see Fig. 1.4),
whose transport matrix is specified by
A = 1, B =d, C = 0, D = 1. (1.3.2)
Figure 1.4: Ray propagation through a drift section.
An analogous result can be obtained for a thin lens system (see Fig. 1.5)
A = 1, B = 0, C =∓ 1
f, D = 1 (1.3.3)
where f is the focal length and the sign specifies the focusing or defocusing
nature of the lens, respectively.
It is now worth noting that ˆO, ˆF, ˆD, indicating the matrices correspond-
ing to a straight section, a focusing and a defocusing lens, respectively (see
Fig. 1.6), we have (see eq. (1.2.4))
ˆO =
( 1 d
0 1
)
= ˆ1 +d ˆσ+ =ed ˆσ+, ˆF =


1 0
− 1
f 1

 = ˆ1− 1
f ˆσ− =e− 1
f ˆσ−,
ˆD =


1 0
1
f 1

 = ˆ1 + 1
f ˆσ− =e
1
f ˆσ− .
(1.3.4)
10We will use the terms free space, drift region or straight section to indicate the same
region.
Matrices, Exponential Operators and Physical Applications 15
Figure 1.5: Ray propagation through a converging lens.
Axis
R1R2
Focal point
Ray propagation
D F O
f
d
Figure 1.6: Symbols for straight section, focusing and defocusing lens, respectively.
The above relations are of noticeable importance since they state that the
propagation through optical systems can be accomplished using the formal-
ism of the Pauli matrices and their algebraic properties. We can derive, from
the previous equation and according to the discussion of the previous sec-
tions, the identity
[
ˆF, ˆO
]
= ˆF ˆO− ˆO ˆF = 2d
f ˆσ3,
[
ˆD, ˆO
]
=−2d
f ˆσ3 (1.3.5)
whose physical meaning is clear: it states that the sequence of the optical
elements, along a transport line, is important. Even though trivial, such a
remark may be helpful in understanding the role of operators in the descrip-
tion of physical devices.
We may ask whether the operator ˆS = eb ˆσ3 (the already introduced
squeeze mapping, see eq. (1.1.22)) can be interpreted as an optical element.
16 Mathematical Methods for Physicists
To this aim, we note that
( r
ϑ
)
0
=eb ˆσ3
( r
ϑ
)
i
=
( ebr
e−bϑ
)
i
. (1.3.6)
We can conclude that it represents the action of a beam expander, whose
practical realization is shown in Fig. 1.7.
Figure 1.7: Beam expander ( Galilean telescope ). We note that the length of the free
region between defocusing and focusing length is f2−f1.
According to the above picture the action of a certain number of succes-
sive optical elements is given by the product of a corresponding number of
matrices. The successive action of a drift section, a thin focusing lens and a
beam expender, can be written as
ˆO ˆF ˆE =ed ˆσ+e− 1
f ˆσ−ebσ 3 (1.3.7)
which has the same structure of the evolution operator discussed in the previ-
ous section and may be exploited to provide a possible physical interpretation
of the characteristic functions. The use of the above few elements about the
operatorial treatment of optical systems allows some interesting conclusions
which are left as exercises.
Exercise 2. a) Prove that the action of two consecutive focusing thin lenses
is equivalent to a single thin lens with focal length
f = f1f2
f1 +f2
. (1.3.8)
b) Consider an optical system composed by a thin lens followed by a straight
section. Prove that such a system is equivalent to the composition
ˆM1 =
(
1 0
− 1
f 1
)( 1 d
0 1
)
=
( 1 d
− 1
f 1− d
f
)
. (1.3.9)
Matrices, Exponential Operators and Physical Applications 17
c) Show that an optical system composed by a drift followed by a thin lens
is equivalent to
ˆM2 =
( 1 d
0 1
)(
1 0
− 1
f 1
)
=
(
1− d
f d
− 1
f 1
)
̸= ˆM1. (1.3.10)
Concerning this result, it is worth stressing that from the mathematical point
of view, it is just a consequence of the non-commutativity of the matrices.
Consider carefully its physical meaning.
d) Show that ˆD ˆO ˆF corresponds to a beam expander with b = ln
(
f1
f2
)
.
e) Discuss the arrangement of ˆD ˆO ˆF and find the condition to obtain a beam
reducer.
f) Finally, develop physical arguments to justify the matrices describing the
optical elements reported in Fig. 1.8 and use the various elements to find the
matrix corresponding to the optical element reported in Fig. 1.3.
Figure 1.8: Optical element and corresponding ABCD matrix.
18 Mathematical Methods for Physicists
We will reconsider the ABCD law of ray beam propagation within the
context of wave optics in the next chapters. The same formalism can further-
more be applied to circuit network like two-port systems of the type shown
in Fig. 1.9.
Figure 1.9: Two port system.
In this case, the ABCD law (represented by the impedances Z) connects
voltages U and currents I and it is straightforward to find a physical justi-
fication for the characterization, reported in Tab. 1.1, of circuit elements in
terms of matrices.
Table 1.1: Circuit element and corresponding ABCD matrix.
Element Matrix Remarks
Series resistor
(
1 −R
0 1
)
R = resistance
Shunt resistor
(
1 0
− 1
R 1
)
R = resistance
Series conductor
(
1 − 1
G
0 1
)
G = conductance
Shunt conductor
(
1 0
−G 1
)
G = conductance
Series inductor
(
1 −Ls
0 1
)
L = inductance ;
s = complex angular frequency
Shunt capacitor
(
1 0
−Cs 1
)
C = capacitance;
s = complex angular frequency
Matrices, Exponential Operators and Physical Applications 19
1.3.2 Quantum Mechanics
In this section, we will discuss the use of 2× 2 matrix operators in quan-
tum mechanical problems like the dynamics of charged particles with spin
interacting with an external magnetic field. The non relativistic Pauli equa-
tion describing the coupling of a spin- 1
2 charged particle with an external
static magnetic field ⃗B writes11
i ℏ∂t Φ =−q ℏ
2m⃗ σ· ⃗B Φ (1.3.11)
where q is the particle charge, m its mass, and Φ is a two-component wave
function. In Fig. 1.10 we have reported a generic spin orientation and if we
Figure 1.10: Total spin direction not aligned with any component(
S = ℏ
√
s(s + 1) =
√
3
2 ℏ
)
.
assume that ⃗B is along the y direction, eq. (1.3.11) can be written as
∂t Φ =ωc (ˆσ+− ˆσ−) Φ, (1.3.12)
whereωc = qB
4m is the cyclotron frequency. According to the discussion of the
previous section, the solution of eq. (1.3.12) is almost straightforward and
we get
Φ = ˆR (ωct) Φ0 , (1.3.13)
where ˆR is the rotation matrix (see eq. (1.1.17)), corresponding to the evo-
lution operator associated with the Hamiltonian of the Schr¨ odinger equation
11The equation has been written in the above form by neglecting the kinetic term and by
retaining the Stern and Gerlach interaction (magnetic moment-magnetic field) only.
20 Mathematical Methods for Physicists
(1.3.12). A more general result can be achieved by casting the above equa-
tion, in the form
∂t Φ = 2iqωc⃗ σ·⃗ nΦ, ⃗ n =
⃗B
B, (1.3.14)
whose solution is again obtained in the form of a rotation.
Let us now consider a two-level quantum system which will be defined
as a quantum system with only two allowed energy states, like an atomic or
molecular system. The two states are coupled by some external field which,
in the example of the Figure 1.11, is an electromagnetic interaction.
Figure 1.11: A schematic representation of an atomic or molecular two-level system cou-
pled by an electromagnetic field.
The Schr¨ odinger equation, accounting for the evolution of such a system, is12
i ℏ∂tΨ = ( ˆH0 + ˆV ) Ψ, (1.3.15)
where V is the interaction potential between upper and lower levels and
ˆH0 is the un-perturbed part of the Hamiltonian, satisfying the eigenvalues
equation
ˆH0ψ± =E±ψ±, E ± =±1
2 ℏω (1.3.16)
with ψ± denoting the wave functions of the two states independent and or-
thogonal13 and E± the relevant energy eigenvalues.
12We assume V independent of time in the sense that its possible variation with time
are very slow compared with the dynamics associated with frequency ωc .
13The concept of orthogonal states will be more carefully discussed in the following.
Here orthogonal means that
∫
ψ∗
+ψ−dV = 0 .
Matrices, Exponential Operators and Physical Applications 21
The evolution of such a system is ruled by the Schr¨ odinger eq. (1.3.15) and
can be written as follows
Ψ =a+(t)ψ+ +a−(t)ψ−. (1.3.17)
By definition, the operator V connects lower and upper levels and does not
contain diagonal elements.
Exercise 3. The equation of motion satisfied by the time dependent coeffi-
cients in eq. (1.3.17) can be inferred from the Schr¨ odinger equation.
Verify that it can be cast in the form
i∂t χ = [ω ˆσz + 2 Ω ˆσy] χ, (1.3.18)
whereχ =
( a+
a−
)
and
i ℏ Ω =⟨ψ+| ˆV|ψ−⟩ =−⟨ψ−| ˆV|ψ+⟩. (1.3.19)
The evolution operator associated with the solution of eq. (1.3.18) is given
by eq. (1.2.22), with
ˆU(t) =e−it⃗Γ·⃗ σ, ⃗Γ≡ (0, 2Ω, ω). (1.3.20)
A straightforward application of eq. (1.2.23) leads to the following time
dependence of the evolution matrix
ˆU(t) =

 cos
(
Γ t
2
)
−iω
sin(Γ t
2)
Γ − 2Ω
Γ sin
(
Γ t
2
)
2Ω
Γ sin
(
Γ t
2
)
cos
(
Γ t
2
)
+iω
sin(Γt
2)
Γ

. (1.3.21)
The probability that the two-level system will be found in the state ψ± is
therefore given by
P±(t) =|⟨ψ±| Ψ⟩|2 =
[
cos
(
Γt
2
)
a± (0)∓ 2Ω
Γ sin(Γt)a∓ (0)
]2
+
+ ω2
Γ2 sin
(
Γt
2
)2
|a± (0)|2,
(1.3.22)
where
|a+ (0)|2 +|a− (0)|2 =|a+ (t)|2 +|a− (t)|2 = 1. (1.3.23)
The examples discussed in the forthcoming sections clarify how the con-
cept of two-level system and the associated formalism apply in a wider con-
text.
22 Mathematical Methods for Physicists
1.3.3 Particle Physics: Kaon Mixing
Before getting into the specific aspects of this problem, let us remind
that, in Particle Physics, the strong interactions are characterized by the
conservation of some quantum numbers like isotopic spin, baryonic number,
strangeness, charm . . . The isotopic spin is a concept introduced to express
the charge independence of the strong interaction, where neutron and proton,
in absence of any other interaction, can be viewed as the manifestation of
the same particle, the nucleon, a kind of two-level system whose upper and
lower states are the proton and the neutron respectively. According to this
picture, we have
N =
( p
n
)
. (1.3.24)
The transitions from the lower to the upper state and vice-versa occur through
an interaction of the type (see Fig. 1.12) 14
ˆH =i
(
g ˆσ+−g∗ˆσ−
)
(1.3.25)
In analogy to the case of two-level atomic or molecular system, the cou-
pling between up and down level is ensured by the exchange of particles, the
mesons, which unlike the photons are massive. Along with the conservation
of the isotopic spin strong interaction are characterized by the conservation
of strangeness, which is an additive quantum number 15. Strangeness S, iso-
topic spin I, baryonic number B and charge Q (normalized to the modulus
of the electron charge) are linked by the relation
Q =I3 + Y
2, Y =B +S, (1.3.26)
where I3 refers to the third component of the isotopic spin and Y is called
hypercharge. The consistency of eq. (1.3.26) can be verified in the case of the
neutron and proton, as shown in Tab. 1.2.
14In the following chapters we will use Feynman diagrams to illustrate processes involv-
ing the interaction between elementary particles. We assume that the reader is familiar
with the relevant qualitative meaning. The chapters are sufficiently self contained and the
reader unfamiliar with the meaning of pittograms appearing in Fig. 1.12 can read them,
before proceeding further.
15Denoting by ˆA an operator such that ˆA|ni⟩ =λi | ni⟩, if|N⟩ =|n1,..., nm⟩ is a
composite system, the operator is said additive if ˆA|N⟩ = (∑m
i=1λi)|N⟩, multiplicative
if ˆA|N⟩ = (∏m
i=1λi)|N⟩ .
Matrices, Exponential Operators and Physical Applications 23
Figure 1.12: Feynman diagram describing the interaction between proton and neutron.
They can be viewed as a two-level system, the coupling between upper and lower level is
ensured by the exchange of a pion. a) This representatiom belongs to an old poit of view,
in b) the pion exchange at the quark level via gluon intermediate exchange is shown (see
below).
Table 1.2: Quantum numbers of the nucleon.
B I3 S Q
p 1 1/2 0 1
n 1 −1/2 0 0
The quarks, the elementary building blocks of strongly interacting parti-
cles, were introduced by Murray Gell-Mann and George Zweig in the early
sixties of the last century and are characterized by the quantum numbers
reported in Tab. 1.3 16
The so called quark model assumes that mesons, are realized as q ¯q com-
bination. The reader can verify that the only allowed combinations are those
reported in Fig. 1.13, which shows the so called octet of pseudo-scalar mesons
(see below). It consists of two iso-doublets of strange mesons (the kaons with
strangeness±1), one iso-triplet of null strangeness mesons (the pions) and
an iso-singlet with strangeness equal to zero: the eta meson (η)17.
16We do not take into account, for the moment c,b , ant t quarks.
17An SU (3) singlet should also be included; the η′.
24 Mathematical Methods for Physicists
Table 1.3: Quarks quantum numbers.
B I3 S Q
u (up) 1/3 1/2 0 2/3
d (down) 1/3 −1/2 0 −1/3
s (strange) 1/3 0 −1 −1/3
Figure 1.13: The I3−S diagram of the pseudo-scalar meson octet.
In Fig. 1.13, we have reported strange and non strange mesons arranged
in aSU (3) octet. The strangeness quantum number S is conserved in strong
interactions but not in weak ones. This means that weak processes do not
distinguish beween the K0 and its antiparticle. From the phenomenologi-
cal point of view, it happens when an initial beam composed by K0, while
propagating, can transform into anti-K0 and then go back to the original con-
figuration. The K0 and anti-K0 states can therefore be mixed. If we view the
two particles as a two-level system, as illustrated in Fig. 1.14, we can envis-
age a kind of oscillation mechanism.
Such a transition can be viewed, at a microscopic scale 18, as a two-level
system in which the lower and upper states are provided by ¯d and s quark
18By microscopic we mean occurring at the level of constituent quarks.
Matrices, Exponential Operators and Physical Applications 25
Figure 1.14: Oscillation of a neutral kaon into its anti-particle.
respectively, coupled by the intermediate charged vector bosonsW (see Chap-
ter 5 for further comments), responsible for the weak interaction decays 19.
If we limit ourselves to the K0− ¯K0 systems, the strangeness operator can
be realized in terms of the ˆt matrix, as it is easily checked. However, being
strangeness not conserved in weak decay processes, it cannot be considered a
good quantum number to study processes involving neutral kaon decays and
to draw useful consequences from the relevant analysis.
Before getting into further details, let us remind a few notions relevant to
weak decay processes which is characterized by the violation of other physical
quantities. The most remarkable is the parity (see Fig. 1.15). We remind
that the operation of parity inversion transforms right handed coordinate
systems into left handed as shown in Fig. 1.15. It can also be viewed as the
Figure 1.15: Parity transformation.
flip in sign of the spatial coordinates, namely
19We are greatly simplifying the discussion. We should indeed include the transitions
due to effects involving heavier quarks ( c,t ) within the context of the so called box dia-
grams.
26 Mathematical Methods for Physicists
ˆP


x
y
z

 =


−x
−y
−z

. (1.3.27)
Particles can be characterized as scalar, vectors, pseudo-scalars or pseudo-
vectors according to their transformation under parity conjugation. By
pseudo-scalar particle, we mean a particle with spin zero but with nega-
tive intrinsic parity, where the parity is a multiplicative quantum number.
Just to give an example from elementary geometry, we note that a vector is a
quantity with negative parity (if reflected it changes sign), the vector product
between two vectors is a pseudo vector, namely a vector with positive parity,
and the scalar product between a pseudo vector and a vector is a pseudo
scalar, namely a scalar with negative parity. In Chapter 10, we will provide
a more extensive treatment of the concepts associated with parity inversion.
As already underlined, parity is not conserved in weak interactions, which
means that mirror symmetry is violated in weak decay process.
In Fig. 1.16 we have reported a
Figure 1.16: Parity violation in beta decay.
sketch of the crucial experiment
which showed the parity violation in
weak decays. It was indeed shown
that in beta decay of 60Co, electrons
are preferentially emitted in the di-
rection opposite to the direction of
the nuclear spin.
A further example of parity is as-
sociated with the C-parity, or char-
ge conjugation, represented by the
operation C| ψ⟩ =| ¯ψ⟩ exchanging
a particle into its anti-particle. The charge conjugation operator applied to
a π+, π− system, with L total angular momentum, yields
C|π+π−⟩ = (−1)L|π−π+⟩. (1.3.28)
We can define the product of the operation of charge and parity conjugations
asCP . An idea of CP transformation applied to the electron is in Fig. 1.17.
Matrices, Exponential Operators and Physical Applications 27
Figure 1.17: CP -conjugation for electrons.
The charge and parity conjugations are not separately conserved in weak
interaction processes. Here we assume that the operation of CP conjugation
is conserved in the neutral kaon oscillation process. It can therefore be con-
sidered a good quantum number to study the decay of neutral kaons. By
noting that
|K0⟩ =| ¯sd⟩, | ¯K0⟩ =| ¯ds ⟩, (1.3.29)
we find
CP |K0⟩ =| ¯ds ⟩ =| ¯K0⟩, C P | ¯K0⟩ =| ¯sd⟩ =|K0⟩,
(1.3.30)
which ensure that they are not eigenvalues of the CP operator, which in
terms of 2× 2 matrix operators can, in this specific case, be written as
CP = ˆh. (1.3.31)
The eigenstates of this operator can be written as a linear combination of
the “pure” K0, ¯K0 states, namely
|KS⟩ = 1√
2
(
|K0⟩ + | ¯K0⟩
)
, |KL⟩ = 1√
2
(
|K0⟩− | ¯K0⟩
)
(1.3.32)
such that
CP |KS⟩ =|KS⟩, CP |KL⟩ =− |KL⟩, (1.3.33)
28 Mathematical Methods for Physicists
where| KS,L ⟩ corresponds to the so called short ( S) and long ( L) lived
K-mesons20. They have different decay modes, and indeed
KS→π+π−, K L→π+π−π0 (1.3.34)
(verify that CP is conserved in the decay modes reported in eq. (1.3.34)).
It is remarkable that these na¨ ıve manipulations have shown that different
combinations of neutral kaon states lead to different particles with different
decay modes. They have different masses and different life-times and we can
further explore their properties by using the simple mathematics developed
so far.
Let us now discuss the evolution of the above two-level system which can
be described through the Hamiltonian
ˆH =


Es−i Γs
2 0
0 EL−i ΓL
2

 (1.3.35)
which is clearly not Hermitian, since it contains two imaginary terms linked to
the decay rates (Γ∝ 1
τ withτ being the life-time of the two particles21). The
evolution of the quantum states (1.3.32) is therefore ruled by the operator
ˆU(t) =

 e−iωst− t
2τs 0
0 e
−iωLt− t
2τL

, ω S,L = ES,L
ℏ , ΓS,L = ℏ
τS,L
.
(1.3.36)
Accordingly, we find
| ΨK(t)⟩ = 1√
2
(
e−iωst− t
2τs |KS⟩ +e
−iωLt− t
2τL |KL⟩
)
. (1.3.37)
The probability that the state will behave as a | ¯K0⟩ is therefore
|⟨ ¯K0| ΨK(t)⟩| 2= 1
4
(
e
− t
τS +e
− t
τL− 2e− t
τ∗ cos ((ωS−ωL)t)
)
, (1.3.38)
20The decay times are given by τS≃ 9· 10−11s, τL≃ 5· 10−8s, hence the name of long
and short living system.
21The system is not closed but open, the particles are not conserved. The non hermitian
part account for this dissipative effect. All open systems in quantum mechanics are treated
with non hermitian operators.
Matrices, Exponential Operators and Physical Applications 29
with τ∗ = 1
2
τSτL
τS +τL
(prove for Exercise).
We want however to point out a different way of computing the previous
probability amplitude. We note indeed that by inverting eq. (1.3.32), we
obtain (
|K0⟩
| ¯K0⟩
)
= 1√
2
(
1 1
1 −1
)(
|KS⟩
|KL⟩
)
. (1.3.39)
The Schr¨ odinger equation describing the evolution of theCP eigenstates
can be transformed into that for the neutral kaon system by simply per-
forming the change of basis given by eq. (1.3.39). In fact, starting from the
Schr¨ odinger equation for the Hamiltonian in eq. (1.3.35),
i∂t
(
|KS⟩
|KL⟩
)
=
(
ωs−i 1
2τS
0
0 ωL−i 1
2τL
)(
|KS⟩
|KL⟩
)
, (1.3.40)
we get, for ˆT = 1√
2
( 1 1
1 −1
)
,
i∂t
(
|K0⟩
| ¯K0⟩
)
= ˆH′
(
|K0⟩
| ¯K0⟩
)
, ˆH′ = ˆT
(
ωs−i 1
2τS
0
0 ωL−i 1
2τL
)
ˆT−1,
(1.3.41)
The effective Hamiltonian yielding the kaon oscillations is therefore
ˆH′ = 1
2


ω+− i
2τ∗
+
ω− + i
2τ∗
−
ω−− i
2τ∗
−
ω+− i
2τ∗
+

, τ ∗
± =τ−1
S ±τ−1
L , ω ± =ωS±ωL.
(1.3.42)
The mechanism ensuring the oscillations is provided by the term ω−,
which also drives the interferential terms in the probability evolution. Since
it is linked to the mass difference of theS,L doublet, it is used in the analysis
of the decay rates to measure the value of this observable 22.
1.4 Cabibbo Angle and See-Saw Mechanism
In this section, we will consider two further examples from Particle Physics
showing that important results can be obtained with a good understanding
22We will not discuss these details, but we remark that mass difference is few eV .
30 Mathematical Methods for Physicists
of phenomenology and with fairly simple mathematical tools. In the previous
section, we have seen that at microscopic level strong or weak process are due
to processes occurring at quark level, as in the classical weak decay process
of the neutron reported in Fig. 1.18.
(a) Neutron decay n→ p +e− + ¯νe at
quark level involved in neutron decay
and beta decay.
(b) Mirror process in nuclear fu-
sion involved in proton-proton
fusion.
Figure 1.18: Weak interaction transformations of u and d quarks. The echanged particle
is the intermediate vector boson W , whose role in weak decay process will be discussed in
Chapter 5.
As already remarked, the kaon oscillation is due to a kind of oscillation in
the s−d quark sector, associated with the strangeness non conservation in
weak interactions.The s−d quarks are not distinguished by these processes
and therefore we can consider a kind of interaction which mixes the s−
d flavours, transforming them into weak interaction eigenstates s′−d′. On
purely phenomenological grounds, s−d quarks can be considered as forming
a two-level system and we can write the interaction Hamiltonian in the form
ˆHs,d = i (χ ˆσ+−χ∗ˆσ−) (1.4.1)
which leads, for χ =χ∗, to the result 23
( d′
s′
)
= eϑˆi
( d
s
)
=
( cosϑ sinϑ
− sinϑ cosϑ
)( d
s
)
, ϑ =χt, (1.4.2)
whereϑ is the Cabibbo angle whose experimental value is close to 13.1 deg
(see Fig. 1.19).
23The angle ϑ is a quantity depending on the strength of the interaction χ and on its
characteristic times. These unknown quantities are all comprised into the phenomenolo-
gical Cabibbo angle.
Matrices, Exponential Operators and Physical Applications 31
Figure 1.19: Cabibbo mixing.
Before discussing a further example taken from Particle Physics, we con-
sider the following Exercise.
Exercise 4. Prove that the following identity is a direct consequence of eq.
(1.2.28)
exp
(
α
( 0 1
1 R
))
= eλ+−eλ−
λ+−λ−
ˆM + λ+eλ−−λ−eλ+
(λ+−λ−)
ˆ1 =
=

 E(α,R ) α sh
(√
∆
2 α
)
α sh
(√
∆
2 α
)
E(α, R) +αR sh
(√
∆
2 α
)

 e
αR
2 ,
(1.4.3)
where
ˆM =α
( 0 1
1 R
)
, λ ± =αR±
√
∆
2 , ∆ = 4 +R2,
E(α, R) = Rα
2 sh
(√
∆
2 α
)
+ cosh
(√
∆
2 α
)
, sh(x) = sinh x
x ,
(1.4.4)
which is a direct consequence of eq. (1.2.28).
32 Mathematical Methods for Physicists
Let us now consider a mass matrix Hamiltonian of the type
ˆm =c2
( 0 m
m M
)
=mc 2
( 0 1
1 R
)
, R = M
m. (1.4.5)
The associated evolution operator is
ˆU(t) = exp
(
−i t
τ
( 0 1
1 R
))
, τ = ℏ
mc 2 (1.4.6)
and it can be explicitly written in a 2 × 2 matrix form, using eq. (1.4.3).
In Figs. 1.20 and 1.21 we have reported the real and imaginary parts of the
diagonal entries of ˆU(t) vs. t, for small and large values of R.
(a) Real (dot line) and Imaginary (con-
tinuous line) parts of ˆU(t)1, 1vs t
τ .
(b) Real (dot line) and Imaginary (con-
tinuous line) parts of ˆU(t)2, 2 vs t
τ .
Figure 1.20: The case R=1.
It is interesting to consider mass matrices of the type (1.4.5) withR>> 1
and explain why they are exploited in Particle Physics within the context of
the so called see-saw mechanism which is invoked to explain why the neutri-
nos have so small masses compared to the other particles
(hint: this exercise is just an invitation to confront himself with an interesting
problem in Particle Physics which we can understand from the mathemati-
cal point of view, but not in its physical aspects, requiring concepts which
will be only marginally covered in this book. The idea beyond the see-saw
mechanism is that the neutrino masses are acquired through an interaction
Matrices, Exponential Operators and Physical Applications 33
(a) Real (dot line) and Imaginary (con-
tinuous line) parts of ˆU(t)1, 1vs t
τ .
(b) Real (dot line) and Imaginary (con-
tinuous line) parts of ˆU(t)2, 2 vs t
τ .
Figure 1.21: The case R=10.
which can be described by means of a mass matrix given by
ˆmν =
( 0 mD
mD MNHL
)
(1.4.7)
where MNHL denotes the mass of a heavy neutral lepton 24, very large com-
pared to the typical mD mass values of the standard model. The coupling
of this lepton to neutrinos determines their masses and their time evolu-
tion characterized by a mixing between different neutrino flavours, known as
neutrino oscillations).
1.5 Gell-Mann and Pauli Matrices
1.5.1 Spin Composition
The rules of composition of angular momentum are well known from
elementary Quantum Mechanics. Two quantum systems with spin- 1
2 each
can be composed to get either a system with spin 0 or 1, according to the
following rule
1
2⊗ 1
2 = 0⊕ 1. (1.5.1)
24The heavy neutral lepton is one of the candidates for the constituents of the dark
matter in the Universe.
34 Mathematical Methods for Physicists
If we denote by|↑⟩ the spin-1
2 states, we can write the previous states in term
of spin wave-functions as
0 → |↑↓⟩−|↓↑⟩√
2 Singletstate,
1 → |↑↑⟩,|↓↓⟩,|↑↓⟩+|↓↑⟩√
2 Tripletstate,
(1.5.2)
where the associated wave functions to spin-1 are different for the different
eigenvalues of sz =−1, 0, 1. The matrices describing spin-1 systems are,
therefore, 3×3 matrices. We invite the reader to derive such a generalization
of the Pauli matrices (for the study of the relevant properties see sec. 5.1).
In the case of the composition of three spin- 1
2 systems we have
1
2⊗ 1
2⊗ 1
2 = 1
2⊕ 1
2⊕ 3
2, (1.5.3)
namely two spin- 1
2 and one spin- 3
2 quantum systems. The state with spin-
3
2 will be specified by the completely symmetric (with respect to the spin
component exchange) wave function
3
2 → |↑↑↑⟩, (1.5.4)
while the states with spin- 1
2 will be specified by the following wave functions
with mixed symmetry 25
1
2 → 1√
6 (2|↓↑↑⟩−|↑↓↑⟩−|↑↑↓⟩ ), 1
2 → 1
2 (|↑↑↓⟩−|↑↓↑⟩ ).
(1.5.5)
It is interesting to derive the 4 × 4 matrices, describing the case of systems
with spin- 3
2 (see sec. 5.1).
Regarding the symmetry properties, the following argument can be used.
Let us first keep in mind that when we combine two spin- 1
2 particles,
the first rule to be respected is the principle of particle indistinguishability .
25It is fairly straightforward to state the symmetry properties associated with the com-
binations of three identical particles. A specific discussion is available in Chapter 5.
Matrices, Exponential Operators and Physical Applications 35
Denoting with ψi,j(1, 2) the wave function, for such a state we must have
|ψi,j(1, 2)|2=|ψi,j(2, 1)|2 (1.5.6)
which yields
ψi,j(1, 2) =eiθψi,j(2, 1) =e2iθψi,j(1, 2). (1.5.7)
Since the double exchange restores the original state, we have e2iθ = 1, for
θ = 0,π . Accordingly, we get that the state may be odd (antisymmetric)
or even (symmetric) under particle permutation, for this reason we find the
following possible configurations (1 =↑, 2 =↓)
Symmetric :



ψ1ψ1,
ψ1ψ2 +ψ2ψ1√
2 ,
ψ2ψ2 ;
Antisymmetric : ψ1ψ2−ψ2ψ1√
2 .
(1.5.8)
The above states corresponds respectively to spin-1 and spin-0. By using the
arrow notation we recover eq. (1.5.2).
Let us now consider the product
( 1
2⊗ 1
2
)
⊗ 1
2 which can be written also as
(0⊕ 1)⊗ 1
2. It is almost straightforward to understand that the symmetric
combination is just provided by
ψ1ψ1ψ1 , ψ1ψ1ψ2 +ψ2ψ1ψ1 +ψ1ψ2ψ1√
3 ,
ψ1ψ2ψ2 +ψ2ψ2ψ1 +ψ2ψ1ψ2√
3 , ψ 2ψ2ψ2 .
(1.5.9)
In terms of “arrow” representation, we end up with the combinations con-
tained in eq. (1.5.5).
Let us now combine the “singlet” ψ0 and the “doublet” ψ1 as
ψ0ψ1−ψ1ψ0→|ψM,2⟩ =|↑↓↑⟩−|↑↑↓⟩√
2 (1.5.10)
36 Mathematical Methods for Physicists
representing a mixed symmetry doublet. The second doublet can be obtained
|ψM,1⟩ = 2|↓↑↑⟩−|↑↑↓⟩−|↑↓↑⟩√
6 (1.5.11)
by exploiting the orthogonal properties of the wave functions, it is indeed
easily checked that
⟨ψM,1|ψM,2⟩ = 0. (1.5.12)
1.5.2 Gell-Mann Matrices
In the previous sections, we have studied the properties of 2×2 Pauli ma-
trices and we have discussed the importance of their application to different
physical problems. We have also seen that 3 × 3 matrices can be exploited
to describe the states of particles with spin 1. This set of three matrices
realizes a representation of SU (2) generators but not a basis for all three
dimensional matrices. A generalization in this sense is provided by a set of 8
matrices known as Gell-Mann matrices which have played a fundamental
role in the development of theSU (3) classification of the elementary particles
and of Quantum Chromo Dynamics.
The Gell-Mann matrices are a direct generalization of the Pauli matrices
and are
ˆλ1 =


0 1 0
1 0 0
0 0 0

, ˆλ2 =


0 −i 0
i 0 0
0 0 0

, ˆλ3 =


1 0 0
0 −1 0
0 0 0

,
ˆλ4 =


0 0 1
0 0 0
1 0 0

, ˆλ5 =


0 0 −i
0 0 0
i 0 0

, ˆλ6 =


0 0 0
0 0 1
0 1 0

,
ˆλ7 =


0 0 0
0 0 −i
0 i 0

, ˆλ8 = 1√
3


1 0 0
0 1 0
0 0 −2

.
(1.5.13)
These matrices satisfy the following relations of commutation
[
ˆλα, ˆλβ
]
=ifαβγ ˆλγ (1.5.14)
where fαβγ plays the same role of the Ricci-Levi Civita tensor in SU (2).
Matrices, Exponential Operators and Physical Applications 37
Exercise 5. Prove that
f123 = 1, f 147 =f165 =f246 =f257 =f345 =f376 = 1
2,
f458 =f678 =
√
3
2 and, whit δαβ Kronecker symbol,
(1.5.15)
Tr
(
ˆλαˆλβ
)
= 2δαβ. (1.5.16)
The physical meaning of the Gell-Mann matrices can be inferred from Fig.
1.22 where we have reported the fundamental quark triplet in the ( I3, Y)
plane. The states on which the ˆλα act are defined by the 3-column vector
Figure 1.22: (I3,Y ) diagram of the fundamental quark triplet.
called the fundamental SU (3) triplet
ϖ =


d
u
s

. (1.5.17)
The first three matrices account for the isotopic spin symmetry (d−u quarks),
the second two for the so calledU-spin symmetry (d−s sector), and the sixth
and seventh account for theu−s sector symmetry (V -spin). The matrix ˆλ8 is
associated with the hypercharge quantum number. It commutes with ˆλ3, as-
sociated with the third component of the isotopic spin, and both hypercharge
and isotopic spin are linked to the charge through the Gell-Mann-Nishima-
Nakano formula (1.3.26).
38 Mathematical Methods for Physicists
1.5.3 Flavors and Spin
In this section, we will see how we can merge the concepts underlying
SU (3) and SU (2) and the quark model. Mesons are constructed by combin-
ing a quark and an anti-quark belonging to the fundamental triplet in all
possible ways. It is not difficult to realize that all the available combinations
are those given in Fig. 1.14 plus a contribution, called SU (3) singlet, whose
wave function in terms of quark content is provided by
|η
′
⟩ = 1√
3
(
|d ¯d⟩+ |u ¯u⟩+|s ¯s⟩
)
. (1.5.18)
Different multiplet wave functions are all mutually orthogonal. In more tech-
nical terms we can say that, by denoting with 3 the fundamental represen-
tation of the SU (3) triplet, mesons belong to the product of representations
3⊗ ¯3 = 1⊕ 8. (1.5.19)
Exercise 6. Prove that the mesons π0, ηare specified by the combinations
|π0⟩ = 1√
2
(
−| d ¯d⟩+|u ¯u⟩
)
, |η⟩ = 1√
6
(
|d ¯d⟩+|u ¯u⟩− 2|s ¯s⟩
)
,
⟨π0| η⟩ = 0
(1.5.20)
(hint: note that we must require ⟨η| η
′
⟩ =⟨π0| η
′
⟩ = 0 ... ).
Baryons are composed of three quarks. The relevant wave functions are
characterized by symmetry properties with respect to the exchange of two
quarks and all the particles with the same symmetry properties belong to
the same SU (3) multiplet.
Exercise 7. Prove that baryons belong to the product of representations
3⊗ 3⊗ 3 = 1⊕ 81⊕ 82⊕ 10. (1.5.21)
The multiplet 10 (see Fig. 1.23a) is completely symmetric, which means
that a particle with quark content u,u,d belonging to this multiplet is char-
acterized by the wave function
| ∆+⟩ = 1√
3 (|uud ⟩+|duu ⟩+|duu ⟩). (1.5.22)
Matrices, Exponential Operators and Physical Applications 39
Any other combination will have mixed symmetry and belong to the multi-
plets with mixed symmetry (namely 8 1,2), where those with the same quark
content as (1.5.22) are (see Chapter 5 for further discussion)
|p1⟩ = 1√
3 (2|duu⟩−| udu ⟩−| uud⟩), (1.5.23)
|p2⟩ = 1√
2 (|uud ⟩−| udu ⟩). (1.5.24)
Eq.(1.5.23) is recognized as the proton (see Fig. 1.23b for the baryon octet
81 where other baryons, along with their quark content, are also shown 26).
(a) Baryon decuplet.
(b) Baryon octect.
Figure 1.23: I3−S diagram.
Finally the singlet 1 has a completely anti-symmetric wave function which
can be easily be realized as (see Chapter 5 for details)
| Γ0⟩ = 1√
6
∑
i,j,k
εijk|ijk ⟩, (1.5.25)
where we made the identification 1→u, 2→d, 3→s. Since the quarks are
fermions with half integer spin, the most correct way to frame their relevant
properties is using a symmetry which includes spin and flavor 27 degrees of
26The second octet contains particles with same quantum numbers of the first. The
interested reader can find more details in the book by D. B. Lichtenberg , Unitary Sym-
metries and Elementary Particles, Benjamin, New York (1978).
27By flavour we mean the SU (3) quantum numbers (isotopic spin and strangeness).
40 Mathematical Methods for Physicists
freedom. From the mathematical point of view this corresponds to studying
the so called direct product SU (3)⊗SU (2).
Proposition 1. The direct product between two groups (G,∗) and (H,◦) is
denoted by G⊗H and it represents a new group in which the set of the
elements is the Cartesian product of the sets of elements of G and H that
is{(g,h ) :g∈G,h∈H} and on these elements we put an operation defined
element wise
(g,h )× (g′,h′) = (g ∗ g′, h◦h′). (1.5.26)
If we denote the fundamental quark triplet by 3 and the fundamental SU (2)
doublet by 1
2, we can write the fundamental representation ofSU (3)⊗SU (2)
as
(
3, 1
2
)
. Mesons will therefore be provided by
(
3, 1
2
)
⊗
(
¯3, 1
2
)
= (1⊕ 8, 0⊕ 1) = (1, 0)⊕ (8, 0)⊕ (1, 1)⊕ (8, 1), (1.5.27)
while for baryons we get
(
3, 1
2
)
⊗
(
3, 1
2
)
⊗
(
3, 1
2
)
=
(
1⊕ 81⊕ 82⊕ 10, 1
2⊕ 3
2
)
. (1.5.28)
Exercise 8. The wave function of mesons and baryons as well should be
written by including the spin so prove that
| ∆++⟩ =|u↑u↑u↑⟩, |π−⟩ = 1√
2 (|d↑u↓⟩−| d↓ ¯u↑⟩ )
(1.5.29)
(hint: recall that ∆ particles have spin- 3
2 and pions spin 0 . . . ).
Furthermore, write the wave functions of vector mesons (mesons with spin- 1
belonging to the multiplet (8, 1)).
Exercise 9. Concerning the proton, we remind that it belongs to the multiplet(
81, 1
2
)
. Prove that its wave function can be written as
|p⟩ = 1
2
√
3
(duu udu uud )


2 −1 −1
−1 2 −1
−1 −1 2




↓↑↑
↑↓↑
↑↑↓

 (1.5.30)
(hint: take into account the symmetry properties of the spin wave-function).
Derive of the wave-functions for all the states in the octet.
Matrices, Exponential Operators and Physical Applications 41
1.5.4 Color and QCD
Since the first proposal of the formulation of the quark model it was
argued, that the SU (3) classification of hadrons, is not fully consistent with
the basic principles of Quantum Mechanics. Baryons, composed by three
quarks, can be accommodated in representations like octets and decuplets.
As already remarked, the baryons belonging to the representation 10 re-
ported in Fig. 1.23a, are symmetric in their quark composition. This fact
implies that states like | Ω−⟩ =| sss ⟩ are not allowed by Pauli principle,
if all three quarks are supposed to be in an S-state. This puzzle is solved if
we assume that the quarks are characterized by a further quantum number,
called color and that they come in three different fundamental colors (red,
blue and green) and that only three different colors may constitute physical
particles, as shown in Fig. 1.24 for the already quoted case of the particle Ω−.
Figure 1.24: Ω− components with three different colors of the constituent quarks.
The use of the adjectives colored and colorless is just a way to indicate new
internal degrees of freedom which label the quarks and may be exploited to
characterize their interaction. The use of red, blue and green colors is jus-
tified by the fact that, in Color Theory, they are called primary colors and
have additive properties like those shown in Fig. 1.25.
From a mathematical point of view, a physical state like Ω− will be characte-
rized by a completely anti-symmetric function in the color variables, namely
| Ω−⟩ = 1√
6
∑
i,j,k
εijk|sisjsk⟩, i,j,k = 1, 2, 3,
i→R, j →G, k →B.
(1.5.31)
42 Mathematical Methods for Physicists
Figure 1.25: Combination of primary additive colors, the overlapping of the three colors
results in a “white” (colorless) image.
The internal structure of a particle like a proton is represented in Fig. 1.26.
The forces which glue the quarks are those carrying the color quantum num-
bers. These forces are therefore mediated by massless particles called gluons,
according to the scheme reported in Fig. 1.27.
Figure 1.26: Internal structure and forces of a proton.
A more thorough discussion is provided in Chapter 5 where we have reported
the (very) elementary tools to deal with Feynman rules concerning differ-
ent types of interactions. The exchange mechanisms are perfectly equivalent
to what happen in the electromagnetic, weak and strong interactions, as sum-
marized in Fig. 1.28. In this picture, the dynamical theory based on the above
picture is called Quantum Chromo Dynamics (QCD) and it is an exact the-
ory, which parallels the QED. The interaction in QED is due to exhange of
photons while inQCD it is carried by eight28 independent color states, which
correspond to the “eight types” or “eight colors” of gluons. Because states
28We do not include the singlet, for the reasons discussed in Chapter 5.
Matrices, Exponential Operators and Physical Applications 43
blue
blue
green
green
green-antiblue
gluon
Figure 1.27: Feynman diagram for an interaction between quarks generated by a gluon
denoted by a curly line.
can be mixed together as discussed above, there are many ways of presenting
these states which are known as theSU (3) color octet (see Tab. 1.4) (verify for
Exercise that they are linked to the Gell-Mann matrices).
Figure 1.28: Fundamental interactions and associated exchange bosons.
We will not add any further comment on QCD. The concepts we have
just mentioned in this section will be sporadically touched on in the next
parts of the book.
1.6 Concluding Remarks
In the previous sections, we have given an introductory idea of how very
na¨ ıve mathematical tools like the evolution operator methods and 2× 2 or
44 Mathematical Methods for Physicists
Table 1.4: SU(3) color octet.
(R ¯B +B ¯R)/
√
2 −i (R ¯B−B ¯R)/
√
2
(R ¯G +G ¯R)/
√
2 −i (R ¯G−G ¯R)/
√
2
(B ¯G−G ¯B )/
√
2 −i (B ¯G−G ¯B )/
√
2
(R ¯R−B ¯B)/
√
2 (R ¯R + B ¯B− 2G ¯G )/
√
6
3× 3 matrices, can be exploited to study different physical problems, which
can be viewed from a common point of view (at least formal).
In this concluding section, we will fill some gaps left open in the previous
discussion by considering, among the other things, the extension of our anal-
ysis to higher order matrices. We will start our discussion with a problem
which can be treated with an extension of the evolution operator method.
1.6.1 Vector Differential Equations and Matrices
Heisenberg equations allow the derivation, from Hamiltonian eq. (1.2.21),
of the spin components motion equations which can be cast in vector form
d
dt⃗ σ=−i
ℏ
[
⃗ σ,ˆH
]
=⃗Ω×⃗ σ, (1.6.1)
which follows from the Pauli matrices commutation properties and from the
following form of vector product
(
⃗ a×⃗b
)
l
=εjkl aj bk . (1.6.2)
Eq. (1.6.1) is formally equivalent to the Euler equation for a rotating body
and it is often used in problems concerning the nuclear magnetic resonance.
From the mathematical point of view, eq. (1.6.1) is an initial value problem
and, according to what we have learned so far, its solution can also be written
applying the formalism of the evolution operator as
⃗ σ(t) =et [⃗Ω]⃗ σ0, (1.6.3)
Matrices, Exponential Operators and Physical Applications 45
where the term in the square brackets denots the operator
[⃗Ω] =⃗Ω×. (1.6.4)
By following the same procedure as before, we expand the exponential oper-
ator and write the solution as
⃗ σ(t) =
∞∑
n=0
tn
n! [⃗Ω]n⃗ σ0 (1.6.5)
which represents a series of iterated vector products. The successive powers
of the operator (1.6.4) should be interpreted as follows
[⃗Ω]0 = ˆ1, [⃗Ω]2 =⃗Ω×(⃗Ω×, [⃗Ω]3 =⃗Ω×(⃗Ω×(⃗Ω .... (1.6.6)
The form of the solution reported in eq. (1.6.5)-(1.6.6) can be useful if trun-
cated after few terms, i.e. if we are interested in solutions valid for short
time intervals, but it is not practical for long time behaviors. The use of the
following property of the internal vector product
⃗ a× (⃗b×⃗ c) = (⃗ a·⃗ c)⃗b− (⃗ a·⃗b)⃗ c (1.6.7)
yields the possibility of writing eq. (1.6.5) in the closed form
⃗ σ(t) = cos(Ωt)⃗ σ0 + sin(Ωt) (⃗ n×⃗ σ0) + [1− cos(Ωt)] (⃗ n·⃗ σ0)⃗ n, (1.6.8)
with ⃗ n=
⃗Ω
Ω. It is called Rodrigues rotation and its geometrical interpre-
tation is reported in Fig. 1.29. The proof follows from the cyclical properties
of the vector product and can be viewed as the result of the extension of the
Euler formula (for further comments see Chapter 6).
The extension of the above method of a vector differential equation so-
lution to the non-homogeneous case is quite straightforward and indeed we
find that the solution of the problem 29
d
dt⃗ σ=⃗Ω×⃗ σ+⃗ γ (1.6.9)
29The general problem of non homogeneous differential equations will be treated in
the following sections. Here we note that, being the non homogeneous term independent
of time, we can solve the problem quite easily by transforming it into a homogeneous
equation. In fact, we have y′ =β−αy, setting η =y− β
α we find η′ =−αη.
46 Mathematical Methods for Physicists
Figure 1.29: Geometrical interpretation of the evolution of ⃗ σ(t). We have assumed for
semplicity ⃗Ω⊥⃗ σ0.
can be written (as the reader can easily prove) in the form
⃗ σ(t) =et [⃗Ω]⃗ σ0 +
∫ t
0
e(t−t′) [⃗Ω]⃗ γdt′ . (1.6.10)
An interesting example of application of the previous result is a genuine
classical problems, relevant to the equation of motion of a body falling under
the influence of the gravity with the inclusion of the Coriolis’ force, namely
d
dt⃗ v=⃗ g− 2⃗Ω×⃗ v. (1.6.11)
The Earth as a rotating frame is shown in Fig. 1.30 in which we have re-
ported the “absolute” frame (XYZ ) and the rotating reference frame ( xyz)
in which the rotation vector ⃗Ω is expressed in terms of the co-latitude angle
λ, while the modulus of the vector is just linked to the Earth angular velocity.
Eq. (1.6.11) is approximate in the sense that it does not contain the cen-
trifugal term [ ⃗Ω]2⃗ rand therefore its validity is limited to first order term
only.
Exercise 10. The use of eqs. (1.6.6) and (1.6.10) yields the following re-
statement of the falling body law
⃗ r(t) =⃗ r0 +⃗ v0 t + 1
2⃗ gt2−⃗Ω×
(
⃗ v0 t2 + 1
3⃗ gt3
)
. (1.6.12)
Matrices, Exponential Operators and Physical Applications 47
Figure 1.30: The fixed and rotating reference frames on the Earth.
Derive the above expression and ponder about its physical meaning.
Find the solution of the equation
d
dt⃗ v=⃗ g− 2⃗Ω×⃗ v−β⃗ v (1.6.13)
which differs from eq. (1.6.11) for the inclusion of a friction contribution.
We have already noted that in spite of the fact that eq. (1.6.1) has been
derived using a quantum formalism, they hold in a classical context too. Such
a statement can be easily verified by deriving the Euler equation of motion
using the Hamilton formalism.
1.6.2 Matrices, Vector Equations and Rotations
By assuming that vector ⃗Ω is specified by components ⃗Ω≡ (Ω1, Ω2, Ω3),
we can cast our equation in the following matrix form (we replace ⃗ σby ⃗S)
d
dt


S1
S2
S3

 =


0 −Ω3 Ω2
Ω3 0 −Ω1
−Ω2 Ω1 0




S1
S2
S3

 (1.6.14)
and the relevant solution as
⃗S = ˆU(t)⃗S0, (1.6.15)
\end{verbatim}
